\documentclass{ieeeaccess}

\usepackage{cite}
\usepackage{amsmath,amssymb,amsfonts}
\usepackage{algorithmic}
\usepackage{graphicx}
\usepackage{textcomp}
\usepackage{booktabs}
\usepackage{tabularx}
\usepackage{array}
\usepackage{url}
\usepackage{balance}

\def\BibTeX{{\rm B\kern-.05em{\sc i\kern-.025em b}\kern-.08em
    T\kern-.1667em\lower.7ex\hbox{E}\kern-.125emX}}

\begin{document}

\history{Date of publication xxxx 00, 0000, date of current version xxxx 00, 0000.}
\doi{10.1109/ACCESS.2026.DOI}

\title{X-MoE: A Task-Aware Mixture-of-Experts Framework for Multi-Class Ethereum Fraud Detection Under Severe Imbalance and Sparse Connectivity}

\author{\uppercase{Nur Haliza Abdul Wahab}\authorrefmark{1,2}}
\address[1]{Digital Trust Centre (DTC), Nanyang Technological University (NTU), Singapore}
\address[2]{Faculty of Computing, Universiti Teknologi Malaysia (UTM), Johor Bahru, Malaysia}

\tfootnote{This work was developed during Visiting Scholar research at the Digital Trust Centre (DTC), Nanyang Technological University (NTU), Singapore.}

\markboth
{Abdul Wahab: X-MoE for Multi-Class Ethereum Fraud Detection}
{Abdul Wahab: X-MoE for Multi-Class Ethereum Fraud Detection}

\corresp{Corresponding author: Nur Haliza Abdul Wahab (e-mail: CS-NUR.HALIZA@assoc.main.ntu.edu.sg).}

\begin{abstract}
Ethereum fraud detection is challenging when labelled malicious addresses are extremely scarce and transaction-graph connectivity is sparse. In the evaluated multi-class BLTE dataset, only 69 Phishing and 99 Scamming addresses are available among 73,033 addresses, while approximately 92.32\% of nodes are isolated in the released graph. This paper proposes X-MoE, a task-aware mixture-of-experts framework for three-class Ethereum address classification. X-MoE combines a conflict-aware task expert with a GraphSAGE expert and performs fraud-class-specific fusion using expert predictions, expert disagreement, and node connectivity. A dedicated Normal-Phishing discrimination objective addresses the extreme imbalance between the dominant Normal class and rare Phishing addresses. Across ten deterministic seeds under a clean topology-visible protocol, X-MoE achieves a Macro-F1 of 0.6575 $\pm$ 0.0712 compared with 0.4916 $\pm$ 0.0336 for GraphSAGE. An eight-variant, ten-seed ablation shows that removing the Normal-Phishing objective causes the only family-wise significant reduction in overall Macro-F1 among the evaluated component removals and substantially degrades Phishing detection. A connectivity-stratified analysis further reveals that removing hard isolate masking improves connected-node discrimination but introduces a small false-positive cost among isolated Normal addresses. Under a strict node-inductive protocol that excludes validation--test connections and direct target--target edges, X-MoE obtains a Macro-F1 of 0.5478 $\pm$ 0.0582 versus 0.4401 $\pm$ 0.0204 for GraphSAGE and wins all ten matched seeds (Holm-adjusted $p = 0.0078$). The findings show that X-MoE provides a robust improvement over conventional GNN baselines within BLTE, while also exposing topology sensitivity and the unequal contribution of individual architectural components.
\end{abstract}

\begin{keywords}
Ethereum, blockchain fraud detection, graph neural networks, mixture of experts, phishing, scamming, class imbalance, graph sparsity, inductive learning.
\end{keywords}

\titlepgskip=-15pt
\maketitle

\section{Introduction}

Ethereum has evolved into a major programmable blockchain infrastructure supporting digital-asset transfers, smart contracts, decentralised applications, and decentralised financial services. Its transparent ledger provides detailed transaction records for security analysis, but its pseudonymous and permissionless nature also enables malicious actors to create and operate addresses with relatively low barriers to entry. Because confirmed blockchain transfers are generally irreversible, reliable identification of fraudulent addresses is important for transaction-risk screening, fraud investigation, and broader blockchain-security applications. Early Ethereum fraud-detection approaches therefore relied heavily on behavioural attributes extracted from account histories, including transaction frequency, transferred value, account lifetime, and transaction-direction statistics. The Benchmark Labeled Transactions Ethereum (BLTE) dataset was introduced to support such analysis using transaction-derived address features and publicly available labels \cite{alemari2021blte}.

Address-level attributes, however, do not capture the complete relational structure of blockchain activity. Ethereum addresses interact through transfers and other transactions, and these interactions can reveal fund-flow patterns, neighbourhood structures, and coordinated behaviours that are not fully represented by local features. This limitation has motivated increasing use of graph-based learning and complementary representations for Ethereum fraud detection. TLMG4Eth integrates transaction-language representations with transaction-similarity and account-interaction graphs \cite{sun2025tlmg4eth}, while KGBERT4Eth combines transaction semantics, a transaction knowledge graph, and expert-curated information \cite{jia2026kgbert4eth}. These studies demonstrate that fraud-relevant evidence can arise from multiple predictive views rather than from a single representation. Nevertheless, the usefulness of graph evidence depends on whether informative relational context is available for a particular address.

This issue becomes particularly pronounced in the three-class BLTE setting considered in this work. A recent extension of BLTE formulates Ethereum address classification using Normal, Phishing, and Scamming labels and provides a controlled benchmark for graph-based multi-class learning \cite{alemari2026multiclass}. The processed dataset evaluated here contains 73,033 Ethereum addresses, comprising 72,865 Normal, 69 Phishing, and 99 Scamming addresses. The two fraud classes together therefore account for less than 0.25\% of the evaluated population. At the same time, approximately 92.32\% of nodes are isolated in the released analysis graph. These characteristics create two interacting difficulties: severe class imbalance limits the amount of supervision available for the minority fraud classes, while extreme graph sparsity prevents most addresses from obtaining neighbourhood evidence through message passing.

\Figure[t!](topskip=0pt, botskip=0pt, midskip=0pt)
[width=\textwidth]
{figures/Figure1.png}
{\parbox[t]{0.86\textwidth}{\raggedright
Motivation for connectivity-aware expert routing in the multi-class BLTE setting. The dataset is severely imbalanced and approximately 92.32\% of nodes are isolated. Connected addresses can exploit both task-level and graph-level evidence, whereas isolated addresses cannot obtain neighbourhood information, motivating adaptive expert routing.
}\label{fig:blte_motivation}}

Fig.~\ref{fig:blte_motivation} illustrates this evidence asymmetry. A connected address can potentially exploit both transaction-derived node attributes and information aggregated from its transaction neighbourhood. An isolated address, in contrast, has no observed neighbours from which a message-passing model can obtain relational evidence and must therefore rely predominantly on its own attributes. The supervision available across classes is similarly uneven: the dominant Normal population provides abundant labelled examples, whereas Phishing and Scamming are represented by only 69 and 99 addresses, respectively. Applying the same predictive pathway uniformly to all nodes may therefore over-rely on unavailable graph context for isolated addresses while failing to preserve sufficiently discriminative decision boundaries for rare fraud classes.

Evidence from the existing multi-class BLTE benchmark reinforces this concern. GraphSAGE achieved the strongest performance among the evaluated GCN, GAT, and GraphSAGE architectures, while a feature-only multilayer perceptron remained statistically competitive after multiple-comparison correction \cite{alemari2026multiclass}. This result suggests that graph information can be useful without being uniformly informative across the dataset. The central problem is therefore not simply to identify a stronger GNN backbone. Instead, a detector should determine when structural information provides useful complementary evidence, when local task evidence should dominate, and whether the relative contribution of these evidence sources should differ across fraud classes.

Existing research addresses parts of this problem from two complementary directions. Ethereum-specific fraud-detection models increasingly combine multiple representations, demonstrating the value of integrating behavioural, semantic, and structural evidence \cite{sun2025tlmg4eth,jia2026kgbert4eth}. In parallel, mixture-of-experts approaches such as CAMERA and MoELLA-GNN demonstrate that specialised experts and adaptive routing can be useful when different predictive mechanisms are informative for different instances \cite{pan2026camera,peng2026moella}. However, these directions have not directly addressed multi-class Ethereum address classification under the combined conditions of extreme minority-class scarcity and highly uneven graph connectivity. In particular, it remains unclear how a task-oriented predictor and a graph-structural predictor should be specialised and combined when relational evidence is unavailable for most nodes, and whether additional routing complexity itself produces better minority-class discrimination.

To address this problem, we propose \textbf{X-MoE}, a task-aware mixture-of-experts framework for classifying Ethereum addresses as Normal, Phishing, or Scamming. X-MoE combines a \textit{Conflict-Aware Multitask Expert}, which emphasises feature- and task-level discrimination, with an independent \textit{GraphSAGE Expert}, which provides structural evidence when useful graph connectivity is available. The framework further incorporates dedicated supervision for the difficult Normal--Phishing distinction and a fraud-class-specific fusion mechanism conditioned on expert outputs, their absolute disagreement, and node connectivity. This formulation enables structural evidence to contribute differently across nodes and output classes instead of imposing uniform reliance on the transaction graph. For isolated nodes, graph contribution is explicitly suppressed, preserving a feature-based prediction pathway. X-MoE is therefore positioned as a problem-specific expert-fusion framework rather than as a claim of generic mixture-of-experts novelty.

The main contributions of this work are as follows:
\begin{itemize}

    \item We propose a task-aware complementary expert framework for multi-class Ethereum fraud detection that integrates a Conflict-Aware Multitask Expert with an independent GraphSAGE Expert. A fraud-class-specific fusion mechanism conditions expert combination on expert predictions, absolute disagreement, and node connectivity, enabling task-level and graph-level evidence to contribute differently across nodes and classes.

    \item We introduce a dedicated Normal--Phishing discrimination objective to strengthen the decision boundary between the dominant Normal class and the extremely rare Phishing class. Controlled ablation shows that removing this objective produces the only family-wise significant reduction in overall Macro-F1 among the evaluated component removals and substantially degrades Phishing discrimination.

    \item We conduct a mechanistic evaluation of expert composition and graph utility through an eight-variant, ten-seed ablation study together with connectivity-stratified analysis. The experiments examine expert composition, routing granularity, specialist supervision, isolate masking, and learned versus fixed fusion, demonstrating that these mechanisms contribute unequally and that additional routing complexity is not uniformly beneficial under BLTE sparsity.

    \item We introduce a strict node-inductive evaluation protocol, denoted STRICT-INDUCTIVE, that excludes validation--test connections and direct target--target edges. Under this setting, X-MoE achieves a Macro-F1 of 0.5478 compared with 0.4401 for GraphSAGE and outperforms GraphSAGE across all ten matched seeds, with a Holm-adjusted paired \(p\)-value of 0.0078. Under the topology-visible clean protocol, X-MoE achieves a mean Macro-F1 of 0.6575 compared with 0.4916 for GraphSAGE. These results show that the comparative advantage of X-MoE persists under stricter node-inductive evaluation within BLTE, without implying temporal or cross-dataset generalisation.

\end{itemize}

The remainder of this paper is organised as follows. Section II reviews related work in Ethereum fraud detection and graph mixture-of-experts learning and positions X-MoE against the closest studies. Section III presents the BLTE dataset, graph characteristics, and formal problem definition. Section IV describes the X-MoE architecture and training objectives. Section V details the experimental protocols, baselines, evaluation metrics, and statistical methodology. Section VI reports the clean, ablation, connectivity, efficiency, and strict-inductive results. Section VII discusses the principal findings and their implications. Section VIII addresses limitations and threats to validity, Section IX summarises reproducibility resources, and Section X concludes the paper.


\section{Background and Related Work}

\subsection{Ethereum Address Fraud Detection and the BLTE Setting}

Ethereum fraud detection has progressively evolved from address-level behavioural modelling toward representations that incorporate relational, semantic, and temporal information. Early approaches commonly represent each address using statistics derived from its transaction history, including transaction frequency, transferred value, account lifetime, and incoming or outgoing activity. The original Benchmark Labeled Transactions Ethereum (BLTE) resource follows this paradigm by providing transaction-derived address features together with publicly available labels for abnormal-address analysis \cite{alemari2021blte}. Such representations remain useful because they provide local behavioural evidence independently of graph connectivity.

Address-level modelling, however, treats an Ethereum account largely through aggregated local behaviour and does not directly represent interactions among addresses. Graph-based approaches address this limitation by modelling accounts as nodes and blockchain interactions as edges. Message-passing GNNs can consequently propagate information across transaction neighbourhoods and encode patterns that are difficult to recover from individual address statistics alone.

The multi-class extension of BLTE is particularly relevant to the present study because it formulates address classification using three labels---Normal, Phishing, and Scamming---and evaluates conventional GCN, GAT, and GraphSAGE architectures under repeated experimental runs \cite{alemari2026multiclass}. That work establishes both the three-class task and the applicability of graph learning to the same underlying setting considered here. Accordingly, neither multi-class Ethereum address classification nor the application of GNNs to this task is claimed as a novelty of X-MoE.

More recent Ethereum fraud-detection studies have increased the expressiveness of graph representations. PDHG constructs a heterogeneous Ethereum transaction graph and applies a heterogeneous graph transformer to preserve differences among account and interaction types for phishing detection \cite{wang2026pdhg}. GTSF incorporates transaction semantics during subgraph construction and applies heterophily-oriented graph learning to distinguish accounts whose interactions do not follow conventional homophily assumptions \cite{song2026gtsf}. Temporal modelling has also been explored; PhishEye represents Ethereum activity as a heterogeneous temporal attributed multi-graph and applies self-supervised temporal graph contrastive learning for phishing detection \cite{wu2026phisheye}.

These studies demonstrate that heterogeneous relations, transaction semantics, temporal information, and specialised graph construction can improve the quality of relational representations when such information is available. They do not, however, remove a fundamental dependence of message-passing models on the availability of neighbourhood evidence. This issue becomes important when graph connectivity varies substantially across target addresses.

\subsection{Graph Learning Under Uneven Connectivity}

A generic message-passing GNN updates the representation of node $v$ by combining its current representation with information aggregated from its neighbourhood. A layer can be expressed abstractly as

\begin{equation}
\mathbf{h}_{v}^{(\ell+1)}
=
\phi^{(\ell)}
\left(
\mathbf{h}_{v}^{(\ell)},
\operatorname{AGG}^{(\ell)}
\left\{
\mathbf{h}_{u}^{(\ell)}
:
u\in\mathcal{N}(v)
\right\}
\right),
\label{eq:generic_message_passing}
\end{equation}

where $\mathbf{h}_{v}^{(\ell)}$ denotes the representation of node $v$ at layer $\ell$, $\mathcal{N}(v)$ is its observed neighbourhood, $\operatorname{AGG}^{(\ell)}(\cdot)$ is a neighbourhood aggregation operator, and $\phi^{(\ell)}(\cdot)$ combines local and aggregated information.

For a connected node, $\mathcal{N}(v)\neq\varnothing$, and the model can exploit both node-level attributes and neighbourhood-derived evidence. For an isolated node,

\begin{equation}
\mathcal{N}(v)=\varnothing,
\label{eq:isolated_neighbourhood}
\end{equation}

so no additional neighbourhood evidence can be obtained through message passing. Although a GNN implementation may still propagate a node's own transformed representation through self-information or residual mechanisms, this is fundamentally different from receiving relational evidence from neighbouring addresses.

This distinction is important in the BLTE setting because connectivity is highly uneven. Improving the expressive power of a GNN---for example through heterogeneous attention, semantic subgraph construction, or temporal encoding---can improve how available neighbourhood information is represented, but it cannot reconstruct relational evidence that is absent from the released graph. Consequently, graph quality and graph availability are separate modelling considerations.

GraphSAGE is particularly relevant to this work because it learns node representations through neighbourhood aggregation while retaining an inductive formulation \cite{hamilton2017graphsage}. In the multi-class BLTE benchmark, GraphSAGE provides the strongest conventional GNN comparison among the evaluated graph architectures \cite{alemari2026multiclass}. Nevertheless, its usefulness remains conditional on the neighbourhood information available to each address. This observation motivates a modelling strategy in which graph-derived predictions complement, rather than universally replace, a task-oriented pathway that can operate from local node evidence.

The implication is that connectivity should be treated as part of the evidence condition of a prediction. A connected address may support both behavioural and structural inference, whereas an isolated address supplies primarily local evidence. X-MoE is developed around this distinction by maintaining complementary task-oriented and graph-structural predictors and allowing their contributions to depend explicitly on node connectivity.

\subsection{Complementary Representations and Expert Routing}

A separate line of Ethereum fraud research has focused on combining multiple predictive views. TLMG4Eth integrates transaction-language representations with transaction-similarity and account-interaction graphs, enabling semantic and structural evidence to be learned jointly \cite{sun2025tlmg4eth}. KGBERT4Eth further combines transaction semantics, a transaction knowledge graph, and expert-curated information to capture contextual, structural, and domain-specific signals \cite{jia2026kgbert4eth}. These studies provide evidence that different representations can encode complementary aspects of fraudulent behaviour.

Self-supervised and multi-view approaches extend this principle further. Meta-IFD introduces generative and contrastive self-supervision for heterogeneous Ethereum interactions and applies the resulting representations to phishing and Ponzi-scheme detection \cite{jin2025metaifd}. LMAE4Eth integrates transaction semantics, masked graph representation learning, and expert knowledge within a multi-view fraud-detection framework \cite{jia2025lmae4eth}. Collectively, these methods show that combining complementary evidence can improve fraud modelling when a single representation is insufficient.

Multi-view representation fusion, however, is conceptually different from mixture-of-experts routing. Multi-view methods generally integrate complementary embeddings or features into a shared predictive representation. A mixture-of-experts model instead preserves specialised predictors and learns how strongly each expert should influence the final decision. This distinction becomes important when different samples do not have access to equivalent evidence.

CAMERA applies expert specialisation to fraud detection by constructing multiple fraud-indicative experts and using contextual information to adapt their contributions in text-attributed graphs \cite{pan2026camera}. Its formulation demonstrates that specialised predictors can capture different fraud-related cues and that adaptive gating can determine how those cues are combined. CAMERA, however, targets unsupervised fraud detection under semantic camouflage in non-blockchain graph settings rather than multi-class Ethereum address classification.

MoELLA-GNN similarly employs multiple graph-oriented experts and gated aggregation in a more general graph-learning setting \cite{peng2026moella}. The framework demonstrates that different graph experts can capture complementary structural signals and that routing can depend on graph context. These studies establish that neither expert specialisation nor adaptive routing is itself novel.

The relevant gap for X-MoE therefore lies not in introducing mixture-of-experts learning as a generic mechanism, but in adapting complementary expert modelling to a setting where two forms of asymmetry occur simultaneously: label evidence is extremely scarce for the minority fraud classes, and graph evidence is unavailable for a large proportion of addresses. Under these conditions, the design question is not simply how to construct a more expressive gate, but how task-level and graph-level predictors should be specialised, supervised, and evaluated when their evidence availability differs across nodes.

\subsection{Positioning of X-MoE}

Table~\ref{tab:xmoe_positioning} positions X-MoE against the studies most closely related to its methodological components. The comparison focuses on the dimensions relevant to the present problem rather than attempting to rank overall model quality.

\begin{table*}[t]
\centering
\caption{Positioning of X-MoE against closely related Ethereum fraud-detection and expert-routing methods}
\label{tab:xmoe_positioning}

\renewcommand{\arraystretch}{1.12}
\setlength{\tabcolsep}{3.2pt}
\scriptsize

\begin{tabularx}{\textwidth}{
>{\raggedright\arraybackslash}p{2.15cm}
>{\raggedright\arraybackslash}p{2.15cm}
>{\raggedright\arraybackslash}p{2.55cm}
>{\centering\arraybackslash}p{1.45cm}
>{\raggedright\arraybackslash}p{2.15cm}
>{\raggedright\arraybackslash}p{2.35cm}
>{\raggedright\arraybackslash}X
}

\toprule

\textbf{Study} &
\textbf{Primary task} &
\textbf{Evidence / representation} &
\textbf{Expert routing} &
\textbf{Connectivity treatment} &
\textbf{Minority-specific supervision} &
\textbf{Evaluation emphasis} \\

\midrule

Al-E'mari et al. \cite{alemari2026multiclass}
&
Normal / Phishing / Scamming
&
Address features + homogeneous transaction graph
&
No
&
Connectivity analysed
&
No dedicated class-pair objective
&
Repeated multi-class GNN comparison
\\[2pt]

TLMG4Eth \cite{sun2025tlmg4eth}
&
Ethereum fraud detection
&
Transaction language + similarity graph + interaction graph
&
No
&
Structural evidence incorporated through multiple graphs
&
Not designed around Normal--Phishing separation
&
Complementary semantic--structural representation
\\[2pt]

KGBERT4Eth \cite{jia2026kgbert4eth}
&
Ethereum fraud detection
&
Transaction semantics + knowledge graph + expert information
&
No
&
Graph structure incorporated through knowledge representation
&
Not designed around Normal--Phishing separation
&
Knowledge-enhanced fraud representation
\\[2pt]

LMAE4Eth \cite{jia2025lmae4eth}
&
Ethereum fraud detection
&
Transaction semantics + masked graph embedding + expert knowledge
&
No
&
Graph representation learned as one component of multi-view fusion
&
No dedicated Normal--Phishing objective
&
Multi-view representation learning
\\[2pt]

CAMERA \cite{pan2026camera}
&
Fraud detection in text-attributed graphs
&
Multiple fraud-indicative experts + contextual information
&
Yes
&
Context influences expert integration
&
Not specific to BLTE fraud classes
&
Adaptive expert specialisation
\\[2pt]

MoELLA-GNN \cite{peng2026moella}
&
General graph learning
&
Multiple graph-oriented experts
&
Yes
&
Graph context influences expert aggregation
&
Not specific to Ethereum minority classes
&
Adaptive graph-expert routing
\\[2pt]

\textbf{X-MoE}
&
\textbf{Normal / Phishing / Scamming}
&
\textbf{Task-oriented + graph-structural experts}
&
\textbf{Yes}
&
\textbf{Explicit degree-aware routing and isolate handling}
&
\textbf{Dedicated Normal--Phishing objective}
&
\textbf{Ablation, connectivity-stratified, and strict node-inductive evaluation}
\\

\bottomrule
\end{tabularx}

\vspace{2pt}

\begin{minipage}{0.97\textwidth}
\scriptsize
\textbf{Note:} The table compares methodological coverage relevant to the present study and is not intended as an overall ranking of model quality. ``Connectivity treatment'' denotes whether graph availability or graph context explicitly affects modelling or analysis; these treatments are not necessarily equivalent across studies.
\end{minipage}

\end{table*}

The comparison highlights three distinctions. First, Ethereum-specific models increasingly exploit richer structural, semantic, and domain-informed representations, but these approaches primarily improve how evidence is represented rather than explicitly maintaining separate task-oriented and graph-oriented predictors. Second, existing mixture-of-experts methods demonstrate the value of expert specialisation and adaptive routing, but they are not designed around the extreme Normal--Phishing--Scamming imbalance and connectivity sparsity of the multi-class BLTE setting. Third, the existing multi-class BLTE benchmark exposes both minority-class difficulty and uneven graph utility, but does not investigate whether complementary predictors should be routed differently according to class and connectivity.

X-MoE is positioned at the intersection of these directions. Its contribution is not the introduction of GNN-based Ethereum classification, multi-view learning, or mixture-of-experts routing in isolation. Instead, it couples a task-oriented expert with an independently parameterised GraphSAGE expert, introduces explicit connectivity-conditioned fusion, and applies targeted supervision to the difficult Normal--Phishing boundary. The resulting architecture is evaluated not only through aggregate predictive performance but also through component ablation, connectivity-stratified analysis, and strict node-inductive testing.

This positioning also defines the claim boundary of the present work. X-MoE does not assume that increasingly complex routing is inherently superior. The experiments explicitly compare class-specific and shared routing, learned and fixed fusion, isolate masking, specialist supervision, and reduced expert configurations. The objective is therefore to determine which mechanisms are empirically supported under the evaluated BLTE setting rather than to attribute any performance improvement indiscriminately to mixture-of-experts complexity.


\section{Problem Formulation and Dataset}

\subsection{Multi-Class BLTE Dataset and Data Provenance}

This study evaluates X-MoE using the processed multi-class Benchmark Labeled Transactions Ethereum (BLTE) setting. The original BLTE resource was developed from labelled Ethereum transaction activity and provides address-level behavioural attributes for abnormal-address analysis \cite{alemari2021blte}. A subsequent study extended the resource to three-class address classification by distinguishing Normal, Phishing, and Scamming addresses and releasing a graph representation suitable for GNN-based evaluation \cite{alemari2026multiclass}.

The experiments in this work operate on the released processed artifacts rather than independently reconstructing the dataset from raw Ethereum blockchain transactions. In particular, the study uses the processed multi-class feature file and graph artifact provided by the upstream project. No additional addresses are introduced, no existing labels are relabelled, and no new fraud labels are assigned during the present study. This distinction is important because the contribution of X-MoE lies in modelling and evaluation rather than in blockchain-data collection or label generation.

The processed artifact contains 73,033 Ethereum addresses, 23 predictive node-level features, and 6,588 directed graph edges. Of the 73,033 addresses, 72,865 are labelled Normal, 69 are labelled Phishing, and 99 are labelled Scamming. The corresponding class proportions are approximately 99.77\%, 0.09\%, and 0.14\%, respectively.

The label ``Normal'' is interpreted according to the available dataset annotation rather than as independently verified evidence of benign behaviour. In this work, a Normal address denotes an address that is not labelled as Phishing or Scamming in the evaluated resource. Public fraud annotations may be incomplete, delayed, or biased toward detected incidents; consequently, the task considered here is supervised classification under the provided BLTE labels rather than definitive adjudication of whether an address is legitimate.

Table~\ref{tab:blte_characteristics} summarises the principal characteristics of the evaluated artifact.

\begin{table*}[t]
\centering
\caption{Characteristics of the evaluated multi-class BLTE artifact}
\label{tab:blte_characteristics}

\renewcommand{\arraystretch}{1.12}
\setlength{\tabcolsep}{5pt}
\scriptsize

\begin{tabularx}{\textwidth}{
>{\raggedright\arraybackslash}p{3.1cm}
>{\centering\arraybackslash}p{2.2cm}
>{\raggedright\arraybackslash}X
}

\toprule
\textbf{Characteristic} &
\textbf{Value} &
\textbf{Modelling implication} \\
\midrule

Total addresses
&
73,033
&
Defines the address-level node-classification population.
\\[2pt]

Predictive node features
&
23
&
Provide local behavioural evidence independently of neighbourhood aggregation.
\\[2pt]

Directed graph edges
&
6,588
&
Provide the released relational structure used by graph-based models.
\\[2pt]

Normal addresses
&
72,865 (99.77\%)
&
Strongly dominant class; overall accuracy can therefore conceal poor minority-class discrimination.
\\[2pt]

Phishing addresses
&
69 (0.09\%)
&
Extremely rare minority class and a central source of majority--minority classification difficulty.
\\[2pt]

Scamming addresses
&
99 (0.14\%)
&
Rare minority class requiring class-specific evaluation.
\\[2pt]

Isolated nodes
&
$\sim$92.32\%
&
Most nodes have no observed neighbours from which a message-passing model can obtain additional relational evidence.
\\[2pt]

Prediction unit
&
Ethereum address
&
The target is address-level class membership rather than transaction-, edge-, or subgraph-level fraud.
\\[2pt]

Output classes
&
Normal / Phishing / Scamming
&
Defines a highly imbalanced three-class classification problem.
\\

\bottomrule
\end{tabularx}

\vspace{2pt}

\begin{minipage}{0.97\textwidth}
\scriptsize
\textbf{Note:} ``Normal'' denotes addresses not labelled as Phishing or Scamming in the evaluated artifact and should not be interpreted as independently verified safe addresses.
\end{minipage}

\end{table*}

\subsection{Input Representation}

Each Ethereum address is represented by a 23-dimensional node-feature vector derived from the processed BLTE artifact. The features describe five complementary aspects of address behaviour: transaction connectivity, transferred amounts, temporal activity, local graph structure, and gas-related behaviour. Table~\ref{tab:blte_features} reports the exact attributes used as model inputs.

The feature matrix is denoted by

\begin{equation}
\mathbf{X}\in\mathbb{R}^{|\mathcal{V}|\times 23},
\label{eq:feature_matrix}
\end{equation}

where each row corresponds to one Ethereum address and each column corresponds to one verified input attribute.

\begin{table*}[t]
\centering
\caption{Node-feature groups in the processed BLTE artifact}
\label{tab:blte_features}

\renewcommand{\arraystretch}{1.15}
\setlength{\tabcolsep}{4pt}
\scriptsize

\begin{tabularx}{\textwidth}{
>{\raggedright\arraybackslash}p{2.45cm}
>{\raggedright\arraybackslash}X
>{\centering\arraybackslash}p{0.9cm}
>{\raggedright\arraybackslash}p{4.0cm}
}

\toprule
\textbf{Feature group} &
\textbf{Exact input attributes} &
\textbf{No.} &
\textbf{Evidence represented} \\
\midrule

Connectivity
&
\texttt{all\_degree},
\texttt{in\_degree},
\texttt{out\_degree},
\texttt{unique in\_degree},
\texttt{unique out\_degree}
&
5
&
Overall transaction connectivity, transaction direction, and diversity of counterparties.
\\[2pt]

Transaction amount
&
\texttt{Avg amount incoming},
\texttt{Avg amount outgoing},
\texttt{Total amount incoming},
\texttt{Total amount outgoing},
\texttt{Max amount incoming},
\texttt{Max amount outgoing},
\texttt{Min amount incoming},
\texttt{Min amount outgoing}
&
8
&
Magnitude and distribution of value transferred into and out of an address.
\\[2pt]

Temporal activity
&
\texttt{Avg time incoming},
\texttt{Avg time outgoing},
\texttt{Total transaction time},
\texttt{Active Duration},
\texttt{Mean time interval},
\texttt{Max time interval},
\texttt{Min time interval}
&
7
&
Timing, activity duration, and regularity of address-level transaction behaviour.
\\[2pt]

Local graph structure
&
\texttt{Clustering coefficient}
&
1
&
Local neighbourhood density and structural cohesion around an address.
\\[2pt]

Gas-related
&
\texttt{Avg gas price},
\texttt{Avg gas limit}
&
2
&
Transaction-cost and resource-allocation behaviour associated with address activity.
\\

\midrule
\textbf{Total}
&
&
\textbf{23}
&
\\
\bottomrule

\end{tabularx}

\vspace{2pt}

\begin{minipage}{0.97\textwidth}
\scriptsize
\textbf{Note:} Feature names are reported exactly as stored in
\texttt{BLTE\_multiclass\_v2.csv}. The fields
\texttt{Address}, \texttt{Scam}, \texttt{attack}, and
\texttt{attack\_encoded} are metadata or label-related fields and are not used as predictive node features.
\end{minipage}

\end{table*}

The connectivity-related attributes describe transaction degree and counterparty diversity. Amount-based attributes characterise the scale and distribution of incoming and outgoing transfers, while temporal attributes capture activity duration and transaction timing. The clustering coefficient represents local structural cohesion, and the gas-related attributes describe transaction-resource behaviour.

These node-level attributes remain available even when a target node has no observed neighbours in the released graph. This distinction is central to X-MoE: local behavioural evidence is available for every address, whereas neighbourhood-derived evidence depends on the observed graph structure.

\subsection{Class and Connectivity Structure}

The evaluated dataset exhibits two strong forms of asymmetry: label imbalance and uneven graph connectivity. Let

\begin{equation}
n_{\mathrm{N}}=72{,}865,\qquad
n_{\mathrm{P}}=69,\qquad
n_{\mathrm{S}}=99
\end{equation}

denote the numbers of Normal, Phishing, and Scamming addresses, respectively. The two fraud classes together therefore represent less than 0.25\% of the evaluated population.

Graph connectivity is similarly uneven. For node $v$, let

\begin{equation}
d_v = d_v^{\mathrm{in}} + d_v^{\mathrm{out}},
\end{equation}

where $d_v^{\mathrm{in}}$ and $d_v^{\mathrm{out}}$ denote the observed in-degree and out-degree in the released graph. We define the connectivity indicator

\begin{equation}
r_v=
\begin{cases}
1, & d_v>0,\\
0, & d_v=0.
\end{cases}
\label{eq:connectivity_indicator}
\end{equation}

Thus, $r_v=0$ identifies an isolated address under the evaluated topology. Approximately 92.32\% of nodes satisfy this condition.

The relevance of this indicator is evidential rather than merely structural. When $r_v=1$, a graph model may potentially combine node attributes with information propagated from neighbouring addresses. When $r_v=0$, no observed neighbourhood exists from which additional relational evidence can be obtained. Graph sparsity is therefore treated as a property of the released BLTE representation used in this work rather than as a universal property of the Ethereum network.

\subsection{Problem Formulation}

Let the processed Ethereum address graph be represented as

\begin{equation}
\mathcal{G}
=
(\mathcal{V},\mathcal{E},\mathbf{X},\mathbf{Y}),
\end{equation}

where $\mathcal{V}$ denotes the set of Ethereum addresses, $\mathcal{E}\subseteq\mathcal{V}\times\mathcal{V}$ is the set of directed graph edges, $\mathbf{X}\in\mathbb{R}^{|\mathcal{V}|\times23}$ is the node-feature matrix, and $\mathbf{Y}$ contains the observed class labels.

Each node $v\in\mathcal{V}$ is assigned one class from

\begin{equation}
\mathcal{C}
=
\{
\mathrm{Normal},
\mathrm{Phishing},
\mathrm{Scamming}
\},
\end{equation}

such that

\begin{equation}
y_v\in\mathcal{C}.
\end{equation}

The objective is to learn a parameterised classifier $f_{\theta}$ that produces a class-probability vector

\begin{equation}
\mathbf{p}_v
=
f_{\theta}
\left(
v,\mathbf{X},\mathcal{E}
\right),
\end{equation}

with

\begin{equation}
\mathbf{p}_v\in[0,1]^3,
\qquad
\sum_{c\in\mathcal{C}}p_{v,c}=1.
\end{equation}

The predicted label is

\begin{equation}
\hat{y}_v
=
\arg\max_{c\in\mathcal{C}}p_{v,c}.
\end{equation}

At a generic level, model learning can be written as

\begin{equation}
\theta^{\ast}
=
\arg\min_{\theta}
\mathcal{L}
\left(
f_{\theta}(\mathcal{G}),
\mathbf{Y}_{\mathrm{tr}}
\right),
\label{eq:generic_objective}
\end{equation}

where $\mathbf{Y}_{\mathrm{tr}}$ denotes the labels available in the training partition and $\mathcal{L}$ denotes the training objective. Section IV specifies the task-aware objectives used by X-MoE.

The key modelling condition is that the evidence available to $f_{\theta}$ differs across nodes. Node-level attributes are available throughout the dataset, whereas neighbourhood evidence is conditional on $r_v$. The learning problem is therefore not merely to maximise three-class classification performance, but to determine how local and relational evidence should contribute when graph availability varies across nodes and class supervision is highly imbalanced.

\subsection{Prediction Assumptions and Scope}

The formulation above is subject to several explicit assumptions.

First, the problem is treated as closed-set classification. Every evaluated address is assumed to belong to one of the three labels provided by the dataset: Normal, Phishing, or Scamming. Previously unseen fraud categories are therefore outside the present prediction space.

Second, node attributes and graph connectivity are defined by the released processed BLTE artifacts. X-MoE does not independently reconstruct historical Ethereum transactions or infer missing graph edges. Absence of an edge is consequently interpreted only with respect to the evaluated graph representation.

Third, the model does not assume that every node has access to equivalent structural evidence. Connectivity is explicitly represented through Eq.~\ref{eq:connectivity_indicator}, and the architecture introduced in Section IV is designed so that graph contribution can be reduced or suppressed when observed neighbourhood information is unavailable.

Fourth, the present formulation does not by itself establish temporal, cross-chain, cross-dataset, or adversarial generalisation. The experiments evaluate held-out address classification within the BLTE resource under multiple graph-visibility protocols. Broader generalisation claims therefore remain outside the scope of the current problem definition.

These assumptions establish the data and prediction conditions under which X-MoE is evaluated. The next section introduces the proposed architecture and explains how task-oriented and graph-structural evidence are combined under these conditions.

\section{Proposed X-MoE Framework}

\subsection{Architecture Overview}

X-MoE is designed for multi-class Ethereum fraud detection under unequal evidence availability. Every address is represented by node-level behavioural attributes, whereas usable neighbourhood information is available only for addresses that are connected in the released graph. The proposed architecture therefore separates feature-based and graph-dependent prediction rather than assuming that graph evidence is uniformly informative for all nodes.

The framework contains two complementary experts. The first is a \textit{Conflict-Aware Multitask Expert}, which contains a feature branch and an internal GraphSAGE branch and performs class-specific feature--graph fusion. The second is an independent \textit{GraphSAGE Expert}, which provides a separately parameterised structural prediction pathway. Their outputs are subsequently combined by a class-specific outer fusion mechanism.

X-MoE consequently contains two levels of routing. The inner gate controls how strongly graph-derived evidence modifies the feature pathway within the Conflict-Aware Multitask Expert. The outer gate controls how strongly the independent GraphSAGE Expert modifies the Conflict-Aware prediction. Both routing stages depend on the availability of graph evidence and are explicitly suppressed for isolated nodes.

This design separates three related questions: whether an address can be classified from its own behavioural attributes, whether its transaction neighbourhood supplies additional class-specific evidence, and how strongly an independent structural prediction should influence the final decision when the two experts disagree.

Figure~\ref{fig:xmoe_architecture} illustrates the overall architecture.

\Figure[t!](topskip=0pt, botskip=0pt, midskip=0pt)
[width=\textwidth]
{figures/figure2.png}
{\parbox[t]{0.82\textwidth}{\raggedright
Proposed X-MoE architecture for multi-class Ethereum fraud detection.
The framework combines a Conflict-Aware Multitask Expert and an independent
GraphSAGE Expert through class-specific, connectivity-aware routing.
The inner gate controls feature--graph integration within the Conflict-Aware
Multitask Expert, while the outer gate combines the two expert predictions
using their logits, disagreement, and node connectivity.
Graph-dependent contributions are suppressed for isolated nodes, providing
a feature-based fallback. Training combines weighted multi-class supervision,
expert-level losses, one-versus-rest objectives, a Scamming-specialist
objective, the Normal--Phishing discrimination objective, and PCGrad.}
\label{fig:xmoe_architecture}}

The input to the framework consists of the 23-dimensional address-level feature vector and the released Ethereum transaction graph. The graph structure is also used to derive the observed degree of each address. The model produces logits for the three classes
Normal, Phishing, and Scamming, followed by a softmax transformation to obtain class probabilities.

Table~\ref{tab:xmoe_components} summarises the principal functional components of the architecture.

\begin{table*}[t]
\centering
\caption{Functional roles of the principal X-MoE components}
\label{tab:xmoe_components}

\renewcommand{\arraystretch}{1.12}
\setlength{\tabcolsep}{5pt}
\scriptsize

\begin{tabularx}{\textwidth}{
>{\raggedright\arraybackslash}p{3.4cm}
>{\raggedright\arraybackslash}p{4.2cm}
>{\raggedright\arraybackslash}X
}
\toprule
\textbf{Component} &
\textbf{Input or mechanism} &
\textbf{Functional role} \\
\midrule

Feature branch
&
Address-level behavioural features
&
Provides a prediction pathway that does not depend on neighbourhood availability.
\\[2pt]

Internal GraphSAGE branch
&
Node features and transaction neighbourhood
&
Encodes local structural evidence within the Conflict-Aware Multitask Expert.
\\[2pt]

Class-specific inner gate
&
Feature representation, graph representation, and node degree
&
Controls the graph contribution independently for Normal, Phishing, and Scamming.
\\[2pt]

Independent GraphSAGE Expert
&
Node features and transaction graph
&
Provides a separately parameterised structural prediction pathway.
\\[2pt]

Class-specific outer gate
&
Expert logits, expert disagreement, and node degree
&
Controls the contribution of the independent structural expert to each output class.
\\[2pt]

Isolate-safe routing
&
Binary connectivity indicator
&
Suppresses graph-dependent contributions when no observed neighbourhood is available.
\\[2pt]

Task-aware objectives
&
Multi-class, expert, one-versus-rest, Scamming, and Normal--Phishing supervision
&
Encourages balanced and class-specific discrimination under severe class imbalance.
\\[2pt]

PCGrad
&
Gradients from multiple objectives
&
Reduces destructive interference among simultaneously optimised objectives.
\\

\bottomrule
\end{tabularx}
\end{table*}


\subsection{Conflict-Aware Multitask Expert}

For each address \(v\), let

\begin{equation}
\mathbf{x}_v \in \mathbb{R}^{d},
\qquad d=23,
\end{equation}

denote its input feature vector.

The Conflict-Aware Multitask Expert first constructs a feature representation using a two-layer fully connected encoder,

\begin{equation}
\mathbf{h}^{F}_v
=
\phi_F(\mathbf{x}_v),
\end{equation}

where \(\phi_F(\cdot)\) denotes the feature encoder. The implementation uses a 64-dimensional hidden representation together with ReLU activation, dropout, and layer normalisation.

In parallel, a graph branch applies two GraphSAGE convolutional layers to the same node features and graph structure,

\begin{equation}
\mathbf{h}^{G}_v
=
\phi_G(\mathbf{X},\mathcal{E})_v.
\end{equation}

GraphSAGE constructs node representations by aggregating local node attributes and neighbourhood information \cite{hamilton2017graphsage}. The graph branch also uses 64-dimensional hidden representations and dropout.

Let

\begin{equation}
d_v
=
d_v^{\mathrm{in}}
+
d_v^{\mathrm{out}}
\end{equation}

denote the observed degree of node \(v\), and define the connectivity indicator

\begin{equation}
c_v
=
\begin{cases}
1, & d_v>0,\\
0, & d_v=0.
\end{cases}
\label{eq:connectivity_indicator}
\end{equation}

The input to the inner routing function is

\begin{equation}
\mathbf{q}^{\mathrm{in}}_v
=
\left[
\mathbf{h}^{F}_v
\parallel
\mathbf{h}^{G}_v
\parallel
\log(1+d_v)
\right].
\end{equation}

The corresponding class-specific graph gate is

\begin{equation}
\boldsymbol{\alpha}_v
=
c_v\,
\sigma
\left(
\mathbf{W}^{\alpha}_2
\,
\mathrm{ReLU}
\left(
\mathbf{W}^{\alpha}_1
\mathbf{q}^{\mathrm{in}}_v
+
\mathbf{b}^{\alpha}_1
\right)
+
\mathbf{b}^{\alpha}_2
\right),
\label{eq:inner_gate}
\end{equation}

where

\begin{equation}
\boldsymbol{\alpha}_v
\in [0,1]^3
\end{equation}

contains separate graph-routing weights for Normal, Phishing, and Scamming.

The feature and graph branches produce class logits

\begin{equation}
\mathbf{z}^{F}_v
=
\mathbf{W}_{F}\mathbf{h}^{F}_v
+
\mathbf{b}_{F},
\end{equation}

and

\begin{equation}
\mathbf{z}^{G}_v
=
\mathbf{W}_{G}\mathbf{h}^{G}_v
+
\mathbf{b}_{G}.
\end{equation}

The Conflict-Aware Multitask Expert combines these logits as

\begin{equation}
\mathbf{z}^{C}_v
=
\mathbf{z}^{F}_v
+
\boldsymbol{\alpha}_v
\odot
\left(
\mathbf{z}^{G}_v
-
\mathbf{z}^{F}_v
\right),
\label{eq:inner_fusion}
\end{equation}

where \(\odot\) denotes element-wise multiplication.

Equation~\ref{eq:inner_fusion} permits graph evidence to contribute differently to each output class rather than enforcing a single feature--graph mixture across the three decisions.


\subsection{Independent GraphSAGE Expert}

The second expert is an independently parameterised two-layer GraphSAGE classifier. For node \(v\), it produces

\begin{equation}
\mathbf{z}^{S}_v
=
E_S(\mathbf{X},\mathcal{E})_v,
\end{equation}

where \(E_S\) denotes the independent GraphSAGE network.

The expert maps the 23 input features to a 64-dimensional hidden representation, applies a second 64-dimensional GraphSAGE layer, and projects the resulting representation to three output logits. ReLU activation and dropout are applied between graph layers.

The role of this expert differs from that of the graph branch inside the Conflict-Aware Multitask Expert. The internal graph branch participates directly in feature--graph interpolation through \(\boldsymbol{\alpha}_v\), whereas the independent GraphSAGE Expert provides a separately trained structural prediction that is subsequently combined at the logit level.

This separation allows the outer routing mechanism to respond not only to the availability of graph structure but also to agreement or disagreement between two independently parameterised prediction pathways.


\subsection{Connectivity-Aware Class-Specific Fusion}

The final prediction is obtained through an outer class-specific gate.

For node \(v\), the outer gate receives

\begin{equation}
\mathbf{q}^{\mathrm{out}}_v
=
\left[
\mathbf{z}^{C}_v
\parallel
\mathbf{z}^{S}_v
\parallel
\left|
\mathbf{z}^{C}_v-\mathbf{z}^{S}_v
\right|
\parallel
\log(1+d_v)
\right].
\label{eq:outer_gate_input}
\end{equation}

The first two terms expose the logits produced by the two experts. Their absolute difference represents prediction disagreement, while the degree term provides explicit information about graph availability.

The class-specific GraphSAGE mixture weights are then

\begin{equation}
\boldsymbol{\beta}_v
=
c_v\,
\sigma
\left(
\mathbf{W}^{\beta}_2
\,
\mathrm{ReLU}
\left(
\mathbf{W}^{\beta}_1
\mathbf{q}^{\mathrm{out}}_v
+
\mathbf{b}^{\beta}_1
\right)
+
\mathbf{b}^{\beta}_2
\right),
\label{eq:outer_gate}
\end{equation}

with

\begin{equation}
\boldsymbol{\beta}_v
\in [0,1]^3.
\end{equation}

The final logits are

\begin{equation}
\mathbf{z}_v
=
\mathbf{z}^{C}_v
+
\boldsymbol{\beta}_v
\odot
\left(
\mathbf{z}^{S}_v
-
\mathbf{z}^{C}_v
\right).
\label{eq:outer_fusion}
\end{equation}

Equivalently, for output class \(k\),

\begin{equation}
z_{v,k}
=
(1-\beta_{v,k})z^{C}_{v,k}
+
\beta_{v,k}z^{S}_{v,k}.
\end{equation}

The model therefore does not impose the same expert mixture on all three classes. Structural evidence may receive a different weight for Normal, Phishing, and Scamming.

The final class probabilities are

\begin{equation}
\mathbf{p}_v
=
\mathrm{softmax}(\mathbf{z}_v),
\end{equation}

and prediction is obtained from

\begin{equation}
\hat{y}_v
=
\arg\max_{k \in \mathcal{C}} p_{v,k}.
\end{equation}


\subsection{Isolate-Safe Routing Property}

The connectivity multiplier in both routing stages gives X-MoE an explicit fallback behaviour for isolated addresses.

When

\begin{equation}
c_v=0,
\end{equation}

the inner routing weights satisfy

\begin{equation}
\boldsymbol{\alpha}_v=\mathbf{0}.
\end{equation}

Equation~\ref{eq:inner_fusion} therefore reduces to

\begin{equation}
\mathbf{z}^{C}_v
=
\mathbf{z}^{F}_v.
\end{equation}

Similarly, the outer gate satisfies

\begin{equation}
\boldsymbol{\beta}_v=\mathbf{0},
\end{equation}

and Equation~\ref{eq:outer_fusion} becomes

\begin{equation}
\mathbf{z}_v
=
\mathbf{z}^{C}_v
=
\mathbf{z}^{F}_v.
\label{eq:isolate_fallback}
\end{equation}

Hence, when no observed neighbourhood is available, the final prediction reduces exactly to the feature pathway.

This property is referred to as \textit{isolate-safe routing}. The term ``safe'' refers only to architectural behaviour: graph-derived logits are prevented from influencing a node that has no observed graph neighbours. It does not imply that the resulting prediction is guaranteed to be correct or risk free.

For connected nodes, both \(\boldsymbol{\alpha}_v\) and \(\boldsymbol{\beta}_v\) remain learnable, allowing structural information to contribute when neighbourhood evidence is present.


\subsection{Task-Aware Training Objectives}

The severe class imbalance and unequal class-specific evidence motivate a set of complementary training objectives.

The primary objective is inverse-frequency weighted multi-class cross-entropy,

\begin{equation}
\mathcal{L}_{\mathrm{main}}
=
-
\sum_{v\in\mathcal{V}_{\mathrm{tr}}}
w_{y_v}
\log p_{v,y_v},
\label{eq:main_loss}
\end{equation}

where \(w_c\) is derived from the frequency of class \(c\) in the training partition.

Two additional weighted cross-entropy objectives supervise the two experts independently,

\begin{equation}
\mathcal{L}^{C}_{\mathrm{exp}}
=
\mathrm{CE}
\left(
\mathbf{z}^{C},
\mathbf{y}
\right),
\end{equation}

and

\begin{equation}
\mathcal{L}^{S}_{\mathrm{exp}}
=
\mathrm{CE}
\left(
\mathbf{z}^{S},
\mathbf{y}
\right).
\end{equation}

Each expert-level objective is assigned weight

\begin{equation}
\lambda_{\mathrm{exp}}=0.25.
\end{equation}

The Conflict-Aware Multitask Expert is further supervised through three class-specific one-versus-rest objectives.

For class \(k\),

\begin{equation}
b_{v,k}
=
\mathbb{I}(y_v=k),
\end{equation}

and the corresponding balanced binary objective is

\begin{equation}
\mathcal{L}^{(k)}_{\mathrm{ovr}}
=
\mathrm{BBCE}
\left(
z^{C}_{v,k},
b_{v,k}
\right),
\end{equation}

where BBCE denotes class-balanced binary cross-entropy. Each one-versus-rest objective receives unit weight.

These objectives provide explicit supervision for Normal, Phishing, and Scamming rather than relying solely on the shared three-class boundary.

A dedicated Scamming-specialist loss is applied to the Scamming output of the independent GraphSAGE Expert,

\begin{equation}
\mathcal{L}_{\mathrm{scam}}
=
\mathrm{BBCE}
\left(
z^{S}_{v,\mathrm{Scamming}},
\mathbb{I}(y_v=\mathrm{Scamming})
\right),
\label{eq:scam_loss}
\end{equation}

with coefficient

\begin{equation}
\lambda_{\mathrm{scam}}=0.5.
\end{equation}


\subsection{Normal--Phishing Discrimination Objective}

The most problem-specific auxiliary objective targets Normal-to-Phishing confusion.

Let

\begin{equation}
\mathcal{V}_{CN}
=
\left\{
v\in\mathcal{V}_{\mathrm{tr}}
:
y_v=\mathrm{Normal},
\;
d_v>0
\right\}
\end{equation}

denote connected Normal training nodes.

For these nodes, X-MoE penalises the probability assigned to the Phishing class,

\begin{equation}
\mathcal{L}_{NP}
=
\frac{1}{|\mathcal{V}_{CN}|}
\sum_{v\in\mathcal{V}_{CN}}
p_{v,\mathrm{Phishing}}.
\label{eq:np_penalty}
\end{equation}

The reference X-MoE configuration uses

\begin{equation}
\lambda_{NP}=0.30.
\label{eq:lambda_np_reference}
\end{equation}

This objective is asymmetric. It is not a separate binary Normal-versus-Phishing classifier. Instead, it discourages the final model from assigning excessive Phishing probability to connected Normal training addresses.

The value in Equation~\ref{eq:lambda_np_reference} is the reference coefficient used in the principal X-MoE configuration. Its contribution is evaluated through ablation, while its sensitivity is examined separately by varying \(\lambda_{NP}\) in the experimental analysis. The reference value is therefore not presented as a uniquely optimal coefficient.


\subsection{Conflict-Aware Optimisation With PCGrad}

The simultaneous training objectives can generate gradients that point in competing directions. X-MoE therefore uses Projected Conflicting Gradient (PCGrad) \cite{yu2020pcgrad} to reduce destructive interference between objectives. PCGrad is adopted as an optimisation mechanism and is not claimed as a methodological contribution of this work.

Let

\begin{equation}
\mathbf{g}_i
=
\nabla_{\theta}\mathcal{L}_i
\end{equation}

and

\begin{equation}
\mathbf{g}_j
=
\nabla_{\theta}\mathcal{L}_j
\end{equation}

denote gradients associated with two objectives.

When

\begin{equation}
\mathbf{g}_i^{\top}\mathbf{g}_j < 0,
\end{equation}

the gradient of objective \(i\) is projected as

\begin{equation}
\mathbf{g}'_i
=
\mathbf{g}_i
-
\frac{
\mathbf{g}_i^{\top}\mathbf{g}_j
}{
\|\mathbf{g}_j\|_2^2
}
\mathbf{g}_j.
\label{eq:pcgrad}
\end{equation}

The resulting projected gradients are subsequently combined to update the shared parameters.

For the reference X-MoE configuration, the optimisation task set contains

\begin{equation}
\left\{
\mathcal{L}_{\mathrm{main}},
\mathcal{L}^{C}_{\mathrm{exp}},
\mathcal{L}^{S}_{\mathrm{exp}},
\mathcal{L}^{(N)}_{\mathrm{ovr}},
\mathcal{L}^{(P)}_{\mathrm{ovr}},
\mathcal{L}^{(S)}_{\mathrm{ovr}},
\mathcal{L}_{\mathrm{scam}},
\mathcal{L}_{NP}
\right\}.
\end{equation}

The corresponding reference coefficients are 0.25 for each expert-level loss, 1.0 for each one-versus-rest objective, 0.5 for the Scamming-specialist objective, and 0.30 for the Normal--Phishing objective.

These values define the reference configuration used in the main comparison. Individual objectives or coefficients are modified only in the explicitly defined ablation and sensitivity experiments.


\subsection{X-MoE Training Procedure}

The complete training process is summarised in Algorithm~1. Feature standardisation, data partitioning, epoch budget, model selection, and the exact experimental protocols are specified separately in Section V.

\begin{figure}[!t]
\centering
\begin{minipage}{0.98\columnwidth}
\small
\textbf{Algorithm 1: Training Procedure of X-MoE}
\vspace{2mm}

\begin{algorithmic}[1]

\STATE \textbf{Input:} node features $\mathbf{X}$, graph $\mathcal{E}$,
training labels $\mathbf{y}$

\STATE Compute node degree $d_v$ and connectivity indicator $c_v$

\FOR{each training epoch}

    \STATE Encode feature representations $\mathbf{h}^{F}$

    \STATE Encode internal GraphSAGE representations $\mathbf{h}^{G}$

    \STATE Compute class-specific inner routing weights
    $\boldsymbol{\alpha}$

    \STATE Obtain Conflict-Aware logits $\mathbf{z}^{C}$

    \STATE Compute independent GraphSAGE logits $\mathbf{z}^{S}$

    \STATE Construct disagreement and connectivity features

    \STATE Compute class-specific outer routing weights
    $\boldsymbol{\beta}$

    \STATE Obtain final logits $\mathbf{z}$ and probabilities $\mathbf{p}$

    \STATE Compute $\mathcal{L}_{\mathrm{main}}$

    \STATE Compute expert-level objectives
    $\mathcal{L}^{C}_{\mathrm{exp}}$ and
    $\mathcal{L}^{S}_{\mathrm{exp}}$

    \STATE Compute three one-versus-rest objectives

    \STATE Compute $\mathcal{L}_{\mathrm{scam}}$

    \STATE Compute $\mathcal{L}_{NP}$ for connected Normal nodes

    \STATE Apply PCGrad across the active optimisation objectives

    \STATE Update model parameters

\ENDFOR

\STATE \textbf{Output:} trained X-MoE parameters

\end{algorithmic}
\end{minipage}
\end{figure}


\subsection{Routing Diagnostics and Interpretation}

The routing architecture exposes several internal quantities that can be inspected directly.

At the inner level, X-MoE produces separate feature--graph routing weights,

\begin{equation}
\boldsymbol{\alpha}_v
=
[
\alpha_{v,N},
\alpha_{v,P},
\alpha_{v,S}
],
\end{equation}

for Normal, Phishing, and Scamming.

At the outer level, the model produces

\begin{equation}
\boldsymbol{\beta}_v
=
[
\beta_{v,N},
\beta_{v,P},
\beta_{v,S}
],
\end{equation}

which indicate the learned contribution of the independent GraphSAGE Expert to each class.

Node degree, expert disagreement, and the changes in predicted probabilities produced by the fusion stages can also be examined.

These quantities provide architectural diagnostics rather than causal explanations. For example, a large \(\beta_{v,P}\) indicates that the final Phishing logit relies more strongly on the independent structural expert for that address. It does not demonstrate that a particular neighbouring address caused the prediction.

The routing variables are therefore interpreted together with controlled ablation and connectivity-stratified evaluation rather than as stand-alone explanations of fraudulent behaviour.

Overall, X-MoE combines feature-based fallback prediction, class-specific structural routing, independent structural expertise, and task-aware supervision within a single multi-class framework. Section V next defines the experimental design used to compare this architecture against feature-only and graph-based baselines, isolate the contribution of its individual mechanisms, test sensitivity to graph visibility, and evaluate the stability of the Normal--Phishing objective.



\section{Experimental Design}

This section defines the experimental protocols used to evaluate X-MoE. The design addresses five complementary questions: overall predictive performance under the released graph representation, the contribution of individual X-MoE mechanisms, behaviour under unequal graph connectivity, robustness to stricter graph visibility, and sensitivity to the Normal--Phishing objective coefficient. Matched data partitions are retained wherever the compared models operate under the same protocol.

\subsection{Evaluation Protocols}

The primary comparison uses a topology-visible clean protocol. For each random seed, the 73,033 BLTE addresses are independently repartitioned using a stratified 70/15/15 train/validation/test split. Stratification preserves the approximate proportions of Normal, Phishing, and Scamming addresses across the three partitions despite the severe class imbalance.

Feature preprocessing is performed independently for every seed. A standard scaler is fitted using only the training-node features and then applied to the corresponding validation and test nodes. Validation and test feature statistics therefore do not enter the preprocessing stage.

The released directed graph topology is retained in the clean experiment. Consequently, although node labels and feature preprocessing follow the train/validation/test partition, addresses assigned to different partitions may remain connected through the released edge set. This setting is therefore referred to as \textit{topology-visible clean evaluation} and should not be interpreted as strict node-inductive evaluation.

All models in the primary comparison are evaluated using the same ten random seeds,

\begin{equation}
\mathcal{S}
=
\{100,101,\ldots,109\}.
\end{equation}

Each model is trained for 300 epochs. Patience-based early stopping is not used. Validation Macro-F1 is monitored throughout training, and the parameter state attaining the highest validation Macro-F1 within the 300-epoch budget is retained for final evaluation on the held-out test partition.

Deterministic execution is enabled for the matched experiments to reduce uncontrolled implementation variation and support paired seed-level statistical analysis.

Table~\ref{tab:evaluation_protocols} summarises the five evaluation settings.

\begin{table*}[t]
\centering
\caption{Summary of the evaluation protocols used in this study}
\label{tab:evaluation_protocols}

\renewcommand{\arraystretch}{1.12}
\setlength{\tabcolsep}{4pt}
\scriptsize

\begin{tabularx}{\textwidth}{
>{\raggedright\arraybackslash}p{2.7cm}
>{\raggedright\arraybackslash}p{2.9cm}
>{\raggedright\arraybackslash}p{3.0cm}
>{\raggedright\arraybackslash}p{3.2cm}
>{\raggedright\arraybackslash}X
}
\toprule
\textbf{Protocol} &
\textbf{Node split} &
\textbf{Feature preprocessing} &
\textbf{Graph visibility} &
\textbf{Primary purpose} \\
\midrule

Topology-visible clean
&
Stratified 70/15/15 split regenerated for each seed
&
Scaler fitted on training nodes only
&
Released topology retained across partitions
&
Primary comparison of MLP, GCN, GAT, GraphSAGE, and X-MoE
\\[2pt]

Ablation
&
Same matched clean splits
&
Same as clean protocol
&
Same released topology
&
Identify the contribution of individual X-MoE mechanisms
\\[2pt]

Connectivity-stratified
&
Same matched clean test sets
&
Same as clean protocol
&
Results separated by observed connectivity
&
Compare behaviour on connected and isolated addresses
\\[2pt]

STRICT-INDUCTIVE
&
Same stratified node partitions
&
Scaler fitted on training nodes only
&
Target--target and cross-evaluation message passing removed
&
Evaluate all five models under stricter topology visibility
\\[2pt]

Hyperparameter sensitivity
&
Same matched clean splits
&
Same as clean protocol
&
Same released topology
&
Evaluate sensitivity to the Normal--Phishing objective coefficient
\\

\bottomrule
\end{tabularx}
\end{table*}


\subsection{Baselines and Model Configuration}

The topology-visible clean comparison contains four baselines and X-MoE. The baselines are selected to distinguish feature-only classification from three established forms of graph representation learning.

The first baseline is a multilayer perceptron (MLP). It receives the same 23-dimensional node-feature vector used by the graph-based models but does not use the transaction edge set. The MLP contains two 64-dimensional hidden layers. Each hidden layer is followed by layer normalisation, ReLU activation, and dropout with probability 0.5 before a final linear layer produces the three class logits.

The MLP is particularly important in the present setting because most nodes in the released BLTE graph are isolated. It therefore provides a direct control for determining how much predictive information is available from the released node attributes without neighbourhood aggregation.

The GCN baseline contains two graph-convolutional layers with 64 hidden units followed by a linear three-class classifier. ReLU activation and dropout are applied between graph layers.

The GAT baseline uses a first graph-attention layer with 32 hidden channels and four attention heads, followed by a second attention layer that produces the final output representation.

The GraphSAGE baseline uses two 64-dimensional neighbourhood-aggregation layers followed by a linear three-class classifier. ReLU activation and dropout are applied between graph layers.

X-MoE follows the architecture defined in Section IV. Its feature and graph representations use 64-dimensional hidden states, dropout is set to 0.5, and the outer class-specific gating network uses a 32-dimensional hidden representation. The reference Normal--Phishing coefficient is

\begin{equation}
\lambda_{NP}=0.30.
\end{equation}

Table~\ref{tab:model_configuration} summarises the principal configuration and experimental role of the five models.

\begin{table*}[t]
\centering
\caption{Principal model configurations used in the topology-visible clean comparison}
\label{tab:model_configuration}

\renewcommand{\arraystretch}{1.12}
\setlength{\tabcolsep}{4pt}
\scriptsize

\begin{tabularx}{\textwidth}{
>{\raggedright\arraybackslash}p{1.9cm}
>{\raggedright\arraybackslash}p{3.0cm}
>{\raggedright\arraybackslash}p{3.2cm}
>{\centering\arraybackslash}p{1.6cm}
>{\raggedright\arraybackslash}X
}
\toprule
\textbf{Model} &
\textbf{Input evidence} &
\textbf{Principal configuration} &
\textbf{Graph} &
\textbf{Experimental role} \\
\midrule

MLP
&
23 node attributes
&
Two 64-dimensional fully connected layers
&
No
&
Feature-only control
\\[2pt]

GCN
&
Node attributes and transaction graph
&
Two 64-dimensional graph-convolution layers
&
Yes
&
Graph-convolution baseline
\\[2pt]

GAT
&
Node attributes and transaction graph
&
32 hidden channels with four attention heads
&
Yes
&
Attention-based graph baseline
\\[2pt]

GraphSAGE
&
Node attributes and transaction graph
&
Two 64-dimensional aggregation layers
&
Yes
&
Neighbourhood-aggregation baseline
\\[2pt]

X-MoE
&
Feature and graph expert pathways
&
64-dimensional expert representations and 32-dimensional outer gate
&
Yes
&
Proposed connectivity-aware mixture-of-experts model
\\

\bottomrule
\end{tabularx}
\end{table*}


\subsection{Implementation and Training Settings}

All compared models are trained under a common implementation protocol. Adam is used as the optimiser with learning rate

\begin{equation}
1\times10^{-3}
\end{equation}

and weight decay

\begin{equation}
5\times10^{-4}.
\end{equation}

The base three-class objective uses inverse-frequency class weighting computed separately from the training partition of each seed. X-MoE additionally uses the task-aware objectives defined in Section IV.

The same processed BLTE artefacts, seed sequence, split-generation procedure, train-only feature standardisation, 300-epoch training budget, checkpoint-selection rule, and metric implementation are used throughout the topology-visible clean comparison.

Table~\ref{tab:implementation_settings} summarises the implementation and training configuration.

\begin{table}[t]
\centering
\caption{Implementation and training settings used in the reported experiments}
\label{tab:implementation_settings}

\renewcommand{\arraystretch}{1.12}
\setlength{\tabcolsep}{5pt}
\scriptsize

\begin{tabular}{ll}
\toprule
\textbf{Setting} & \textbf{Value} \\
\midrule

Python & 3.11.9 \\
PyTorch & 2.14.1+cpu \\
PyTorch Geometric & 2.8.0.post1 \\
Execution device & CPU \\
Optimizer & Adam \\
Learning rate & $1\times10^{-3}$ \\
Weight decay & $5\times10^{-4}$ \\
Training epochs & 300 \\
Random seeds & 100--109 \\
Hidden dimension & 64 \\
Dropout & 0.5 \\
X-MoE outer-gate hidden dimension & 32 \\
Reference $\lambda_{NP}$ & 0.30 \\
Model selection & Highest validation Macro-F1 \\

\bottomrule
\end{tabular}
\end{table}

The software versions reported in Table 7 correspond to the verified CPU environment used for the reported experiments.

\subsection{Evaluation Metrics and Statistical Analysis}

Macro-F1 is used as the primary evaluation metric because it assigns equal importance to Normal, Phishing, and Scamming irrespective of class prevalence. For the class set \(\mathcal{C}\),

\begin{equation}
F_{1}^{\mathrm{macro}}
=
\frac{1}{|\mathcal{C}|}
\sum_{c\in\mathcal{C}}
F_{1,c}.
\label{eq:macro_f1}
\end{equation}

This metric is particularly important for BLTE because a model can achieve high overall accuracy by favouring the dominant Normal class while providing weak discrimination for the minority fraud classes.

Accuracy is reported as a supporting global measure. Class-specific precision, recall, and F1 are calculated for Normal, Phishing, and Scamming. Average precision (AP) is additionally reported for Phishing and Scamming because it evaluates minority-class ranking quality across decision thresholds and remains informative under severe class imbalance.

Results are reported as mean and standard deviation across the ten matched seeds.

The clean- and strict-protocol baseline comparisons use the two-sided exact Wilcoxon signed-rank test applied to matched seed-level observations. Holm correction is applied separately across four X-MoE-versus-baseline comparisons within each protocol. Clean-to-strict Macro-F1 changes form a separate family of five within-model comparisons. The strict minority comparison against GraphSAGE forms another family of four metrics: Phishing F1/AP and Scamming F1/AP. For an all-zero vector of paired differences, as observed for MLP clean-to-strict Macro-F1, the reported p-value is set to one; this records no observed change and is not a general equivalence claim. The paired design ensures that two configurations being compared are evaluated using identical seed-specific train, validation, and test partitions.

For the ablation study, the primary analysis uses an exact two-sided paired sign-flip permutation test on the mean seed-level difference between each variant and A0. The test enumerates all $2^{10}=1{,}024$ sign assignments and compares the absolute permuted mean difference with the observed absolute mean difference. Its exactness is conditional on joint sign exchangeability under the null hypothesis. Holm correction is applied across the seven A0-versus-ablation comparisons. Ablation effect sizes use paired Cohen's $d_z$, with differences defined as variant minus A0; 95\% confidence intervals for the mean differences use 100{,}000 paired bootstrap resamples.

Where multiple related hypotheses are evaluated within the same comparison family, the resulting \(p\)-values are adjusted using the Holm procedure to control the family-wise error rate. Statistical significance is assessed at

\begin{equation}
\alpha=0.05.
\end{equation}

Practical effect magnitude is evaluated together with statistical significance. For paired observations, the standardised mean difference is

\begin{equation}
d_z
=
\frac{\overline{\Delta}}
{s_{\Delta}},
\label{eq:dz}
\end{equation}

where \(\overline{\Delta}\) is the mean paired seed-level difference and \(s_{\Delta}\) is the standard deviation of the paired differences.

Where appropriate, 95\% confidence intervals are additionally reported for seed-level estimates or paired differences. Statistical interpretation therefore considers mean performance, variability, paired significance, multiplicity correction, confidence intervals, and effect magnitude rather than relying on mean performance alone.


\subsection{Ablation Design}

The ablation study evaluates whether the observed behaviour of X-MoE can be attributed to specific architectural and training mechanisms rather than simply to increased model capacity.

Eight configurations are evaluated using the same ten matched topology-visible clean splits and the same 300-epoch training budget.

Table~\ref{tab:ablation_design} defines the ablation variants.

\begin{table*}[t]
\centering
\caption{X-MoE ablation configurations}
\label{tab:ablation_design}

\renewcommand{\arraystretch}{1.12}
\setlength{\tabcolsep}{5pt}
\scriptsize

\begin{tabularx}{\textwidth}{
>{\centering\arraybackslash}p{0.9cm}
>{\raggedright\arraybackslash}p{4.2cm}
>{\raggedright\arraybackslash}X
}
\toprule
\textbf{ID} &
\textbf{Configuration} &
\textbf{Purpose} \\
\midrule

A0
&
Full X-MoE
&
Reference configuration containing both experts, class-specific routing, isolate masking, and all task-aware objectives.
\\[2pt]

A1
&
Conflict-Aware Multitask Expert only
&
Removes the independent GraphSAGE Expert to assess the contribution of the outer structural pathway.
\\[2pt]

A2
&
GraphSAGE-expert configuration
&
Evaluates a simpler graph-expert configuration within the task-aware training framework.
\\[2pt]

A3
&
Shared routing
&
Replaces class-specific routing with shared routing to test whether class-conditioned mixture weights are necessary.
\\[2pt]

A4
&
No isolate masking
&
Allows graph-dependent contributions for isolated addresses and evaluates the role of isolate-safe routing.
\\[2pt]

A5
&
No Scamming-specialist objective
&
Removes the dedicated Scamming auxiliary objective.
\\[2pt]

A6
&
No Normal--Phishing objective
&
Removes the dedicated Normal--Phishing discrimination objective.
\\[2pt]

A7
&
Fixed 50:50 outer mixture
&
Replaces learned outer expert routing with equal expert contributions.
\\

\bottomrule
\end{tabularx}
\end{table*}

The full configuration is not assumed a priori to outperform every reduced variant. The purpose of the ablation study is to determine which architectural and optimisation mechanisms receive empirical support and which components introduce only limited benefit or measurable trade-offs.

Matched seed-level Macro-F1 is the primary ablation measure. Class-specific F1 and AP are additionally examined when the removed component is explicitly associated with a minority-class objective.

Particular attention is given to A6 because it directly evaluates the Normal--Phishing objective introduced in Section IV.


\subsection{Connectivity-Stratified Evaluation}

The connectivity-stratified analysis evaluates whether model behaviour changes according to the availability of graph neighbourhood evidence.

Each test node is assigned to one of two subsets,

\begin{equation}
\mathcal{V}_{\mathrm{conn}}
=
\{v:d_v>0\},
\end{equation}

and

\begin{equation}
\mathcal{V}_{\mathrm{iso}}
=
\{v:d_v=0\}.
\end{equation}

The connected subset contains addresses with at least one observed graph neighbour, whereas the isolated subset contains addresses for which no neighbourhood evidence is available in the released graph representation.

In the evaluated BLTE splits, all Phishing and Scamming test addresses belong to the connected subset, whereas the isolated test subset contains only Normal addresses. The two subsets therefore support different interpretations.

For connected nodes, Macro-F1 together with Phishing and Scamming F1 and AP are examined to assess minority-class discrimination where graph information is available.

For isolated nodes, fraud recall cannot be estimated because no isolated fraud-labelled test addresses are present. The relevant quantity is therefore the false-positive rate among isolated Normal addresses,

\begin{equation}
\mathrm{FPR}_{\mathrm{iso}}
=
\frac{
\#\{\text{isolated Normal nodes predicted as fraud}\}
}{
\#\{\text{isolated Normal nodes}\}
}.
\label{eq:fpr_iso}
\end{equation}

The principal connectivity comparison is between the full X-MoE configuration A0 and the no-isolate-masking variant A4. This comparison directly evaluates the trade-off introduced by isolate-safe routing: whether suppressing graph contribution protects isolated Normal addresses and whether that protection alters fraud discrimination among connected nodes.

No conclusion about the detection of isolated Phishing or isolated Scamming addresses is drawn because such labelled test cases are absent from the evaluated representation.


\subsection{Strict Node-Inductive Evaluation}

The topology-visible clean protocol controls node partitioning and feature preprocessing but retains the released graph topology across partitions. To determine whether the comparative advantage of X-MoE depends on evaluation-time topology visibility, a second protocol, denoted STRICT-INDUCTIVE, imposes stronger restrictions on graph construction.

During training, only training nodes and edges whose two endpoints belong to the training partition are visible,

\begin{equation}
\mathcal{G}_{\mathrm{train}}
=
(
\mathcal{V}_{\mathrm{train}},
\mathcal{E}_{\mathrm{train}\leftrightarrow\mathrm{train}}
).
\end{equation}

During validation, the visible graph contains training and validation nodes. Only training--training and training--validation edges are retained. Validation--validation edges are excluded, and test nodes are not visible.

During testing, the visible graph contains training and test nodes. Training--training and training--test edges are retained, whereas test--test edges are removed and validation nodes are excluded.

The original direction of each retained edge is preserved.

Consequently, a validation or test address may receive relational information through training nodes but cannot aggregate information from another address belonging to the same evaluation partition. The protocol therefore removes target--target message passing and prevents cross-evaluation topology visibility.

The protocol is described as strict \textit{node-inductive} evaluation rather than temporal, cross-chain, or cross-dataset generalisation. All evaluated addresses remain drawn from the same BLTE resource.

GCN, GAT, GraphSAGE, and X-MoE are re-evaluated under STRICT-INDUCTIVE using the same ten matched seeds and the same 300-epoch budget.

All five models, including the feature-only MLP, are evaluated under STRICT-INDUCTIVE using ten matched seeds. MLP does not use graph edges, and its per-seed Macro-F1 is identical to the clean result for every seed. It provides an attribute-only reference for interpreting changes in the graph-dependent models.


\subsection{Hyperparameter Sensitivity}

The ablation study determines whether the dedicated Normal--Phishing objective contributes to X-MoE when it is present or removed. A separate sensitivity experiment evaluates whether model performance depends strongly on the particular coefficient assigned to that objective.

The coefficient is varied over

\begin{equation}
\lambda_{NP}
\in
\{0,\;0.10,\;0.30,\;0.50,\;1.00\}.
\label{eq:lambda_np_sensitivity}
\end{equation}

All other architectural, optimisation, preprocessing, loss-weight, split, and checkpoint-selection settings are held fixed. Each value is evaluated using the same ten matched clean-protocol seeds and the same 300-epoch training budget.

The value

\begin{equation}
\lambda_{NP}=0.30
\end{equation}

is retained as the reference configuration used in the principal X-MoE experiments.

The sensitivity analysis is not used to retrospectively select the coefficient that produces the numerically highest test Macro-F1. Instead, it evaluates whether the reference value lies within a broader region of stable performance.

Macro-F1 is the primary sensitivity measure. Phishing F1 and Phishing AP are examined because the objective directly targets Normal-to-Phishing confusion. Scamming F1 and Scamming AP are monitored simultaneously to identify possible trade-offs with the second minority fraud class. Accuracy is retained as a supporting measure only.

The sensitivity analysis is restricted to \(\lambda_{NP}\) because this coefficient corresponds directly to one of the problem-specific mechanisms introduced in Section IV. Broader test-set-driven searches over hidden dimensions or gate widths are not used to select the reported architecture.

Together, the five evaluation settings provide complementary evidence. The topology-visible clean experiment compares X-MoE with both feature-only and graph-based baselines. The ablation study identifies which X-MoE mechanisms contribute to its behaviour. The connectivity-stratified analysis examines the effect of unequal neighbourhood availability. STRICT-INDUCTIVE evaluates all five models under more restrictive topology visibility, with MLP providing an attribute-only reference. Finally, the sensitivity experiment tests whether the Normal--Phishing mechanism depends on a narrowly tuned coefficient.




\section{Experimental Results}

This section reports the empirical results obtained under the evaluation protocols defined in Section V. The analysis progresses from the primary topology-visible clean comparison to paired statistical testing, class-specific behaviour, component ablation, connectivity-stratified analysis, strict node-inductive evaluation, hyperparameter sensitivity, and computational cost. This organisation separates overall predictive performance from the mechanisms and evidence conditions that produce it.

\subsection{Overall Performance Under the Clean Protocol}

Table~\ref{tab:clean_overall_results} compares X-MoE with the feature-only MLP and the three conventional GNN baselines under the topology-visible clean protocol.

\begin{table}[t]
\centering
\caption{Overall performance under the topology-visible clean protocol across ten seeds}
\label{tab:clean_overall_results}

\renewcommand{\arraystretch}{1.12}
\setlength{\tabcolsep}{4pt}
\scriptsize

\begin{tabular}{lcc}
\toprule
\textbf{Model} &
\textbf{Accuracy} &
\textbf{Macro-F1} \\
\midrule

GCN &
$0.8811 \pm 0.0329$ &
$0.3637 \pm 0.0205$ \\

GAT &
$0.5917 \pm 0.0336$ &
$0.2665 \pm 0.0171$ \\

MLP &
$0.9672 \pm 0.0081$ &
$0.3715 \pm 0.0085$ \\

GraphSAGE &
$0.9882 \pm 0.0021$ &
$0.4916 \pm 0.0336$ \\

\textbf{X-MoE} &
$\mathbf{0.9969 \pm 0.0012}$ &
$\mathbf{0.6575 \pm 0.0712}$ \\

\bottomrule
\end{tabular}
\end{table}

X-MoE achieves the highest mean Macro-F1 of 0.6575. GraphSAGE is the strongest conventional graph baseline with 0.4916, followed by the feature-only MLP at 0.3715, GCN at 0.3637, and GAT at 0.2665.

Relative to GraphSAGE, X-MoE improves mean Macro-F1 by

\begin{equation}
0.6575-0.4916=0.1659,
\end{equation}

corresponding to an absolute gain of 16.59 percentage points and a relative improvement of approximately 33.74\%.

The MLP result provides an important control. Its Macro-F1 slightly exceeds that of GCN, demonstrating that the 23 released address-level attributes contain useful predictive information even without message passing. Nevertheless, GraphSAGE substantially exceeds the feature-only MLP, and X-MoE improves further beyond GraphSAGE. The main result therefore cannot be explained solely by the predictive strength of the local node attributes.

Accuracy should be interpreted cautiously because approximately 99.77\% of the evaluated addresses are labelled Normal. High accuracy can therefore coexist with poor minority-class discrimination. Macro-F1 provides the more informative primary comparison because it weights the three classes equally.

X-MoE also exhibits greater Macro-F1 variation across seeds than the conventional baselines. This variation is important in the present dataset because each split contains only a small number of Phishing and Scamming examples. The following analyses therefore examine both paired seed-level consistency and class-specific behaviour.


\subsection{Statistical Comparison With Baselines}

Because all clean-protocol models use the same ten seed-specific partitions, their Macro-F1 values can be compared using paired statistical analysis.

Table~\ref{tab:clean_statistical_results} reports the mean paired Macro-F1 difference, 95\% confidence interval of the paired difference, Holm-adjusted exact Wilcoxon signed-rank \(p\)-value, paired standardised effect size \(d_z\), and the number of seeds in which X-MoE achieves the higher Macro-F1.

\begin{table*}[t]
\centering
\caption{Paired statistical comparison of X-MoE with the clean-protocol baselines}
\label{tab:clean_statistical_results}

\renewcommand{\arraystretch}{1.12}
\setlength{\tabcolsep}{5pt}
\scriptsize

\begin{tabular}{lccccc}
\toprule
\textbf{Comparison} &
\(\boldsymbol{\Delta}\) \textbf{Macro-F1} &
\textbf{95\% CI} &
\textbf{Holm-adjusted \(p\)} &
\(\boldsymbol{d_z}\) &
\textbf{X-MoE wins} \\
\midrule

X-MoE vs. GCN &
0.2938 &
[0.2450, 0.3426] &
0.007812 &
4.309 &
10/10 \\

X-MoE vs. GAT &
0.3910 &
[0.3382, 0.4438] &
0.007812 &
5.299 &
10/10 \\

X-MoE vs. MLP &
0.2860 &
[0.2387, 0.3333] &
0.007812 &
4.323 &
10/10 \\

X-MoE vs. GraphSAGE &
0.1659 &
[0.1216, 0.2101] &
0.007812 &
2.681 &
10/10 \\

\bottomrule
\end{tabular}
\end{table*}

X-MoE achieves a higher Macro-F1 than every baseline in all ten matched seeds. The exact two-sided Wilcoxon signed-rank test gives an unadjusted \(p\)-value of 0.001953 for each comparison. After Holm correction across the four baseline comparisons, all remain significant with an adjusted \(p\)-value of 0.007812.

GraphSAGE provides the most conservative comparison because it is the strongest conventional graph baseline. The mean paired X-MoE--GraphSAGE difference is 0.1659, with a 95\% confidence interval of [0.1216, 0.2101] and \(d_z=2.681\).

The feature-only comparison leads to a similar conclusion. X-MoE exceeds MLP by a mean paired difference of 0.2860, with a 95\% confidence interval of [0.2387, 0.3333] and \(d_z=4.323\). Thus, the X-MoE advantage cannot be attributed solely to information already present in the node attributes.

Together, the paired results show that the clean-protocol ranking is not driven only by differences in ten-seed means. The X-MoE advantage is consistently positive across all evaluated partitions and remains statistically significant after family-wise correction.


\subsection{Minority-Class Performance}

The aggregate Macro-F1 improvement is examined further at the class level because BLTE is extremely imbalanced.

Table~\ref{tab:clean_class_results} reports class-specific F1 for all five clean-protocol models.

\begin{table}[t]
\centering
\caption{Class-specific F1 under the topology-visible clean protocol}
\label{tab:clean_class_results}

\renewcommand{\arraystretch}{1.12}
\setlength{\tabcolsep}{4pt}
\scriptsize

\begin{tabular}{lccc}
\toprule
\textbf{Model} &
\textbf{Normal} &
\textbf{Phishing} &
\textbf{Scamming} \\
\midrule

GCN &
0.9367 &
0.0132 &
0.1412 \\

GAT &
0.7430 &
0.0036 &
0.0528 \\

MLP &
0.9837 &
0.0434 &
0.0874 \\

GraphSAGE &
0.9944 &
0.1023 &
0.3782 \\

\textbf{X-MoE} &
\textbf{0.9988} &
\textbf{0.5374} &
\textbf{0.4363} \\

\bottomrule
\end{tabular}
\end{table}

All stronger models achieve high F1 on the dominant Normal class. The more consequential differences occur for the minority classes.

The largest improvement is observed for Phishing. X-MoE increases mean Phishing F1 from 0.1023 for GraphSAGE to 0.5374 and increases Phishing AP from 0.1871 to 0.5098.

The corresponding precision--recall behaviour shows how this improvement is obtained. GraphSAGE achieves Phishing precision of 0.0564 and recall of 0.5609, whereas X-MoE achieves precision of 0.5966 and recall of 0.5045. Thus, X-MoE substantially improves the precision of Phishing predictions while retaining similar recall.

For Scamming, the improvement is smaller. Scamming F1 increases from 0.3782 for GraphSAGE to 0.4363 for X-MoE, while average precision increases from 0.3667 to 0.5070. X-MoE also increases Scamming precision from 0.2755 to 0.3214 and recall from 0.6867 to 0.7600.

Fig.~\ref{fig:minority_performance} summarises the minority-class comparison between GraphSAGE and X-MoE.

\Figure[t!](topskip=0pt, botskip=0pt, midskip=0pt)
[width=\textwidth]
{figures/figure3.png}
{\parbox[t]{0.82\textwidth}{\raggedright
Minority-class performance of GraphSAGE and X-MoE under the topology-visible clean protocol across ten matched seeds. X-MoE substantially improves both F1 and average precision for Phishing, while smaller but positive improvements are observed for Scamming. Values represent mean scores across the ten evaluated seeds.}
\label{fig:minority_performance}}

The class-specific results show that the increase in overall Macro-F1 is not driven primarily by the already dominant Normal class. Instead, the strongest improvement occurs for Phishing, motivating the component-level analysis of the Normal--Phishing objective.


\subsection{Ablation Analysis}

Table~\ref{tab:ablation_results} reports the mean Macro-F1 for the eight X-MoE ablation configurations.

\begin{table*}[t]
\centering
\caption{Ablation results under the topology-visible clean protocol}
\label{tab:ablation_results}

\renewcommand{\arraystretch}{1.12}
\setlength{\tabcolsep}{5pt}
\scriptsize

\begin{tabularx}{\textwidth}{
>{\centering\arraybackslash}p{0.8cm}
>{\raggedright\arraybackslash}p{4.8cm}
>{\centering\arraybackslash}p{2.0cm}
>{\centering\arraybackslash}p{2.0cm}
>{\raggedright\arraybackslash}X
}
\toprule
\textbf{ID} &
\textbf{Configuration} &
\textbf{Macro-F1} &
\textbf{$\Delta$ vs. A0} &
\textbf{Interpretation} \\
\midrule

A0 &
Full X-MoE &
0.6575 &
-- &
Reference configuration
\\[2pt]

A1 &
Conflict-Aware Multitask Expert only &
0.6525 &
$-0.0050$ &
Small aggregate change after removing the independent outer expert
\\[2pt]

A2 &
GraphSAGE-expert configuration &
0.6558 &
$-0.0017$ &
Retains nearly all A0 Macro-F1 with substantially lower complexity
\\[2pt]

A3 &
Shared routing &
0.6917 &
$+0.0342$ &
Class-specific routing is not required to maximise aggregate Macro-F1
\\[2pt]

A4 &
No isolate masking &
\textbf{0.7071} &
$\mathbf{+0.0496}$ &
Improves aggregate performance but alters isolated-node behaviour
\\[2pt]

A5 &
No Scamming-specialist objective &
0.6557 &
$-0.0018$ &
Only a small aggregate change
\\[2pt]

A6 &
No Normal--Phishing objective &
0.5526 &
$-0.1049$ &
Largest degradation among the evaluated ablations
\\[2pt]

A7 &
Fixed 50:50 outer mixture &
0.6912 &
$+0.0337$ &
Learned outer routing is not uniformly superior to fixed mixing
\\

\bottomrule
\end{tabularx}
\end{table*}

The ablation study does not support the claim that every component of the complete architecture independently increases aggregate Macro-F1. A3, A4, and A7 achieve higher mean Macro-F1 than A0.

The strongest negative effect occurs in A6. Removing the Normal--Phishing objective decreases mean Macro-F1 from 0.6575 to 0.5526,

\begin{equation}
\Delta F_{1}^{\mathrm{macro}}=-0.1049.
\end{equation}

Across the ten matched seeds, A6 does not outperform A0. The exact two-sided paired sign-flip permutation test gives an unadjusted \(p\)-value of 0.001953, a Holm-adjusted \(p\)-value of 0.013672, and

\begin{equation}
d_z=-2.284.
\end{equation}

Among the overall ablation comparisons, A6 is the only variant that remains statistically significant after family-wise correction.

The degradation is concentrated primarily in Phishing. Removing the Normal--Phishing objective decreases Phishing F1 from 0.5374 to 0.2379,

\begin{equation}
\Delta F_{1,\mathrm{Phishing}}
=
-0.2995,
\end{equation}

and reduces Phishing AP from 0.5098 to 0.3357,

\begin{equation}
\Delta AP_{\mathrm{Phishing}}
=
-0.1741.
\end{equation}

These findings provide the clearest mechanism-level evidence in the ablation analysis. The targeted objective materially improves the difficult Normal--Phishing decision boundary under the evaluated BLTE setting.

The other ablations qualify the interpretation of architectural complexity. Removing the independent GraphSAGE expert in A1 or the Scamming-specialist objective in A5 produces only small aggregate changes. A2 is especially notable because it achieves a Macro-F1 of 0.6558, only 0.0017 below A0, despite using a substantially simpler configuration.

The empirical contribution of X-MoE should therefore not be described simply as a consequence of adding more routing or expert capacity. Task-aware supervision receives stronger direct empirical support than several of the additional routing mechanisms.


\subsection{Connectivity-Stratified Results}

The higher aggregate Macro-F1 of A4 cannot be interpreted adequately without considering node connectivity. Table~\ref{tab:connectivity_results} therefore compares A0 and A4 after separating connected nodes from isolated Normal nodes.

\begin{table*}[t]
\centering
\caption{Effect of removing isolate masking under connectivity-stratified evaluation}
\label{tab:connectivity_results}

\renewcommand{\arraystretch}{1.12}
\setlength{\tabcolsep}{5pt}
\scriptsize

\begin{tabular}{lccc}
\toprule
\textbf{Metric} &
\textbf{A0: Full X-MoE} &
\textbf{A4: No Isolate Masking} &
\textbf{Difference} \\
\midrule

Connected-node Macro-F1 &
0.6524 &
\textbf{0.7273} &
$+0.0749$ \\

Connected Phishing F1 &
0.5374 &
\textbf{0.5576} &
$+0.0202$ \\

Connected Phishing AP &
0.5107 &
\textbf{0.5149} &
$+0.0042$ \\

Connected Scamming F1 &
0.4363 &
\textbf{0.6301} &
$+0.1938$ \\

Connected Scamming AP &
0.5070 &
\textbf{0.5862} &
$+0.0791$ \\

Isolated Normal FPR &
\textbf{0.000000} &
0.000307 &
$+0.000307$ \\

\bottomrule
\end{tabular}
\end{table*}

Removing isolate masking increases connected-node Macro-F1 from 0.6524 to 0.7273. The largest class-specific change occurs for Scamming F1, which increases from 0.4363 to 0.6301.

The improvement is accompanied by a small false-positive cost among isolated Normal nodes. A0 produces an isolated Normal false-positive rate of zero, whereas A4 increases this rate to 0.000307. This corresponds to approximately 3.1 additional false-positive predictions per 10,000 isolated Normal addresses.

The result therefore represents a trade-off rather than a simple superiority relation. Removing isolate masking improves discrimination where relational evidence is present, particularly for Scamming, whereas the reference isolate-safe design enforces more conservative behaviour when no observed neighbourhood exists.

No conclusion can be drawn about detection of isolated Phishing or Scamming addresses because all fraud-labelled test nodes in the evaluated splits are connected.


\subsection{Performance Under STRICT-INDUCTIVE}

Table~\ref{tab:strict_overall_results} reports Macro-F1 under STRICT-INDUCTIVE.

\begin{table}[t]
\centering
\caption{Macro-F1 under STRICT-INDUCTIVE across ten matched seeds}
\label{tab:strict_overall_results}

\renewcommand{\arraystretch}{1.12}
\setlength{\tabcolsep}{5pt}
\scriptsize

\begin{tabular}{lc}
\toprule
\textbf{Model} &
\textbf{Macro-F1} \\
\midrule

MLP &
$0.3715 \pm 0.0085$ \\

GCN &
$0.3583 \pm 0.0206$ \\

GAT &
$0.2771 \pm 0.0236$ \\

GraphSAGE &
$0.4401 \pm 0.0204$ \\

\textbf{X-MoE} &
$\mathbf{0.5478 \pm 0.0582}$ \\

\bottomrule
\end{tabular}
\end{table}

X-MoE achieves the highest mean Macro-F1 among all five evaluated models under the stricter protocol. It exceeds MLP by 0.1763, GCN by 0.1895, GAT by 0.2707, and GraphSAGE by 0.1076.

X-MoE achieves the higher Macro-F1 in all ten matched seeds against each of the four baselines. All four comparisons remain significant after Holm correction with an adjusted \(p\)-value of 0.0078125.

The minority-class comparison against GraphSAGE is reported in Table~\ref{tab:strict_minority_results}.

\begin{table}[t]
\centering
\caption{Minority-class comparison under STRICT-INDUCTIVE}
\label{tab:strict_minority_results}

\renewcommand{\arraystretch}{1.12}
\setlength{\tabcolsep}{4pt}
\scriptsize

\begin{tabular}{lcc}
\toprule
\textbf{Metric} &
\textbf{GraphSAGE} &
\textbf{X-MoE} \\
\midrule

Phishing F1 &
0.0825 &
\textbf{0.3611} \\

Phishing AP &
0.1765 &
\textbf{0.3617} \\

Scamming F1 &
0.2472 &
\textbf{0.2855} \\

Scamming AP &
0.2452 &
\textbf{0.3371} \\

\bottomrule
\end{tabular}
\end{table}

For Phishing, X-MoE increases F1 from 0.0825 to 0.3611 and AP from 0.1765 to 0.3617. Both differences remain significant after Holm correction across the four minority metrics, with adjusted \(p=0.0078125\) for each.

The Scamming result is more nuanced. F1 increases from 0.2472 to 0.2855, but the difference is not statistically significant after correction, with a Holm-adjusted \(p\)-value of 0.275391. In contrast, Scamming AP increases from 0.2452 to 0.3371 and remains significant with an adjusted \(p\)-value of 0.027344.

The strict evaluation therefore provides stronger evidence for improved Phishing discrimination than for threshold-dependent Scamming classification. Importantly, the experiment evaluates node-inductive behaviour within BLTE and should not be interpreted as evidence of temporal, cross-chain, or cross-dataset generalisation.


\subsection{Clean-to-Strict Performance Change}

Fig.~\ref{fig:clean_strict_macro_f1} compares the topology-visible clean protocol with STRICT-INDUCTIVE.

\Figure[t!](topskip=0pt, botskip=0pt, midskip=0pt)
[width=\textwidth]
{figures/figure4_clean_vs_strict.png}
{\parbox[t]{0.82\textwidth}{\raggedright
Comparison of all five models under the topology-visible clean protocol and STRICT-INDUCTIVE across ten matched seeds. MLP Macro-F1 is unchanged between protocols. Error bars indicate standard deviation. X-MoE achieves the highest mean Macro-F1 under both protocols, although its performance decreases when target-side graph visibility is restricted. The X-MoE--GraphSAGE gap decreases from 0.1659 under the clean protocol to 0.1076 under STRICT-INDUCTIVE, corresponding to 64.9\% descriptive gap retention.}
\label{fig:clean_strict_macro_f1}}

Table~\ref{tab:clean_strict_change} quantifies the corresponding changes.

\begin{table*}[t]
\centering
\caption{Change in Macro-F1 from the topology-visible clean protocol to STRICT-INDUCTIVE}
\label{tab:clean_strict_change}

\renewcommand{\arraystretch}{1.12}
\setlength{\tabcolsep}{6pt}
\scriptsize

\begin{tabular}{lcccc}
\toprule
\textbf{Model} &
\textbf{Clean} &
\textbf{Strict} &
\textbf{Relative change} &
\textbf{Holm-adjusted \(p\)} \\
\midrule

MLP &
0.3715 &
0.3715 &
$0.00\%$ &
1.000000 \\

GCN &
0.3637 &
0.3583 &
$-1.47\%$ &
0.644531 \\

GAT &
0.2665 &
0.2771 &
$+3.96\%$ &
0.058594 \\

GraphSAGE &
0.4916 &
0.4401 &
$-10.47\%$ &
0.009766 \\

X-MoE &
0.6575 &
0.5478 &
$-16.69\%$ &
0.015625 \\

\bottomrule
\end{tabular}
\end{table*}

X-MoE exhibits the largest relative decline, decreasing by 16.69\%, and GraphSAGE decreases by 10.47\%; both changes remain significant after Holm correction across five models. MLP Macro-F1 is unchanged at every seed. GCN changes only slightly, and the small GAT increase is not significant after correction (adjusted \(p=0.058594\)).

These changes should be interpreted as sensitivity to available graph evidence rather than as direct evidence of leakage in the clean protocol. STRICT-INDUCTIVE alters the relations available during message passing, and X-MoE is explicitly designed to exploit relational information when it is available.

Despite the decrease, the X-MoE--GraphSAGE advantage remains substantial. The gap changes from

\begin{equation}
0.6575-0.4916=0.1659
\end{equation}

under the clean protocol to

\begin{equation}
0.5478-0.4401=0.1076
\end{equation}

under STRICT-INDUCTIVE.

The strict-setting gap is approximately 64.9\% of the clean-setting gap. This percentage is descriptive and is not treated as a separate statistical test.


\subsection{Sensitivity to the Normal--Phishing Coefficient}

Table~\ref{tab:lambda_sensitivity_results} evaluates the sensitivity of X-MoE to the coefficient of the Normal--Phishing objective.

\begin{table*}[t]
\centering
\caption{Sensitivity of X-MoE to the Normal--Phishing coefficient across ten seeds}
\label{tab:lambda_sensitivity_results}

\renewcommand{\arraystretch}{1.12}
\setlength{\tabcolsep}{5pt}
\scriptsize

\begin{tabular}{ccccc}
\toprule
\(\boldsymbol{\lambda_{NP}}\) &
\textbf{Accuracy} &
\textbf{Macro-F1} &
\textbf{Phishing F1} &
\textbf{Scamming F1} \\
\midrule

0.00 &
$0.9939 \pm 0.0018$ &
$0.5526 \pm 0.0451$ &
$0.2379 \pm 0.0766$ &
$0.4226 \pm 0.1082$ \\

0.10 &
$0.9964 \pm 0.0014$ &
$0.6489 \pm 0.0662$ &
$0.5438 \pm 0.1363$ &
$0.4044 \pm 0.1007$ \\

\textbf{0.30} &
$\mathbf{0.9969 \pm 0.0012}$ &
$0.6575 \pm 0.0712$ &
$0.5374 \pm 0.1433$ &
$\mathbf{0.4363 \pm 0.1010}$ \\

0.50 &
$0.9965 \pm 0.0019$ &
$\mathbf{0.6580 \pm 0.0746}$ &
$\mathbf{0.5551 \pm 0.1647}$ &
$0.4203 \pm 0.1174$ \\

1.00 &
$0.9965 \pm 0.0019$ &
$0.6523 \pm 0.0776$ &
$0.5402 \pm 0.1569$ &
$0.4180 \pm 0.1255$ \\

\bottomrule
\end{tabular}
\end{table*}

The largest sensitivity occurs when the objective is removed entirely. Setting

\begin{equation}
\lambda_{NP}=0
\end{equation}

reduces mean Macro-F1 to 0.5526 and Phishing F1 to 0.2379.

Once the objective is active, performance remains comparatively stable. Across

\begin{equation}
\lambda_{NP}\in[0.10,1.00],
\end{equation}

mean Macro-F1 ranges only from 0.6489 to 0.6580.

Although \(\lambda_{NP}=0.50\) produces the numerically highest mean Macro-F1, its difference from the reference value 0.30 is only

\begin{equation}
0.6580-0.6575\approx0.0005.
\end{equation}

This difference is negligible relative to the observed seed-level variation. The sensitivity analysis is therefore not used to retrospectively redefine the reference coefficient.

The reference value \(\lambda_{NP}=0.30\) lies within a broad high-performing region rather than representing an isolated optimum. Together with the A6 ablation, this result indicates that the presence of the targeted Normal--Phishing objective is substantially more important than fine adjustment of its coefficient.


\subsection{Model Complexity and Computational Cost}

The ablation experiments also reveal a trade-off between predictive performance and architectural complexity.

Table~\ref{tab:efficiency_results} reports the trainable parameter count and observed training time for the X-MoE configurations.

\begin{table*}[t]
\centering
\caption{Model complexity and observed training time for the X-MoE ablation configurations}
\label{tab:efficiency_results}

\renewcommand{\arraystretch}{1.12}
\setlength{\tabcolsep}{5pt}
\scriptsize

\begin{tabular}{lccc}
\toprule
\textbf{Configuration} &
\textbf{Parameters} &
\textbf{Training time (s)} &
\textbf{Macro-F1} \\
\midrule

A0 Full X-MoE &
33,775 &
239.39 &
0.6575 \\

A1 Conflict-Aware only &
21,865 &
131.96 &
0.6525 \\

A2 GraphSAGE-expert configuration &
11,459 &
66.80 &
0.6558 \\

A3 Shared routing &
33,775 &
238.56 &
0.6917 \\

A4 No isolate masking &
33,775 &
226.92 &
0.7071 \\

A5 No Scamming specialist &
33,775 &
221.22 &
0.6557 \\

A6 No Normal--Phishing objective &
33,775 &
202.35 &
0.5526 \\

A7 Fixed 50:50 outer mixture &
33,324 &
220.00 &
0.6912 \\

\bottomrule
\end{tabular}
\end{table*}

The full X-MoE configuration contains 33,775 trainable parameters and requires approximately 239.39 seconds per reported training run in the evaluated CPU environment.

A2 provides the clearest efficiency comparison. It contains 11,459 parameters, approximately 34\% of the parameter count of A0, and requires approximately 28\% of its observed training time while obtaining a Macro-F1 of 0.6558, only 0.0017 below the reference configuration.

The results therefore do not support a simple claim that the complete architecture is necessary for maximum predictive efficiency. Instead, they reveal a performance--complexity--behaviour trade-off. The full model retains explicit class-specific routing and isolate-safe behaviour, whereas simpler variants can provide similar or even higher aggregate Macro-F1 with lower or comparable computational cost.

Overall, the experiments provide five principal findings. First, X-MoE significantly improves Macro-F1 over MLP, GCN, GAT, and GraphSAGE under the topology-visible clean protocol and wins all ten matched seeds against each baseline. Second, the largest class-specific improvement occurs for Phishing rather than the already dominant Normal class. Third, the Normal--Phishing objective provides the strongest mechanism-level evidence, with its removal causing the largest and only family-wise significant aggregate degradation among the evaluated ablations. Fourth, isolate-safe routing introduces a measurable trade-off between stronger connected-node discrimination and conservative behaviour on isolated Normal addresses. Fifth, X-MoE retains a statistically significant advantage under STRICT-INDUCTIVE, although its absolute performance decreases when relational evidence is more tightly restricted.



\section{Discussion}

The experimental results show that X-MoE improves balanced multi-class Ethereum fraud detection under the evaluated BLTE setting, but they also reveal that the source of this improvement is more specific than a generic architectural-complexity effect. The strongest evidence concerns minority-class discrimination, targeted supervision at the Normal--Phishing boundary, and the conditional value of graph information. At the same time, several simplified routing variants perform competitively or even exceed the reference configuration in aggregate Macro-F1. The resulting interpretation is therefore one of complementary evidence and task-aware supervision rather than universal superiority of a more complex mixture-of-experts design.

\subsection{Why X-MoE Improves Minority-Class Detection}

The main advantage of X-MoE lies in improving minority-class discrimination rather than further increasing performance on the dominant Normal class. Under the topology-visible clean protocol, GraphSAGE and X-MoE both achieve near-perfect Normal F1, whereas the largest separation occurs for Phishing. X-MoE increases Phishing F1 from 0.1023 to 0.5374 and Phishing AP from 0.1871 to 0.5098 relative to GraphSAGE.

The feature-only MLP provides an additional perspective on this result. Its Macro-F1 of 0.3715 slightly exceeds that of GCN, indicating that the 23 address-level behavioural attributes already contain useful predictive information. However, MLP remains well below GraphSAGE at 0.4916 and X-MoE at 0.6575. The improvement obtained by X-MoE therefore cannot be explained by node attributes alone.

This pattern supports the use of complementary predictive pathways. The feature-oriented pathway can exploit behavioural attributes even when an address has no observed neighbours, whereas the GraphSAGE pathway can exploit relational information when graph connectivity is available. The value of X-MoE does not arise simply from replacing one representation with another, but from preserving both sources of evidence and allowing them to contribute differently across classes and nodes.

The Phishing result is especially informative because only 69 labelled Phishing addresses are available in the evaluated dataset. Learning a reliable decision boundary from such limited positive evidence is difficult when the overwhelming majority of training samples are Normal. X-MoE addresses this problem through several forms of task-aware supervision: inverse-frequency weighting reduces majority-class dominance, one-versus-rest objectives reinforce class-specific discrimination, and the dedicated Normal--Phishing objective explicitly penalises a failure mode associated with the rare Phishing class.

The combined evidence therefore suggests that the principal predictive gain should not be attributed simply to adding another GNN branch. A more defensible interpretation is that complementary feature and structural evidence become more effective when combined with supervision targeted at the most difficult majority--minority decision boundary.


\subsection{What the Ablation and Sensitivity Results Reveal}

The ablation study challenges a simple interpretation in which every component of the full X-MoE architecture is necessary for strong predictive performance. Several reduced configurations match or exceed the reference model in aggregate Macro-F1.

A2 is particularly informative. The GraphSAGE-expert configuration attains a Macro-F1 of 0.6558, only 0.0017 below the full model, while using substantially fewer trainable parameters and less training time. Shared routing in A3 and fixed outer fusion in A7 also exceed A0 in mean Macro-F1, while removing isolate masking in A4 produces the highest aggregate Macro-F1 among the evaluated variants.

These results indicate that additional routing flexibility does not automatically translate into higher predictive performance. Class-specific gates provide a structured mechanism for analysing how evidence is combined, but the BLTE setting does not provide evidence that every routing component is individually necessary for maximising Macro-F1.

The strongest ablation evidence instead supports the dedicated Normal--Phishing objective. Removing this objective in A6 reduces Macro-F1 from 0.6575 to 0.5526 and produces the only family-wise significant degradation in aggregate Macro-F1 among the evaluated component removals. The corresponding Phishing F1 decreases from 0.5374 to 0.2379, showing that the loss of performance is concentrated largely at the class boundary the objective was designed to address.

The hyperparameter-sensitivity experiment strengthens this interpretation. When

\begin{equation}
\lambda_{NP}=0,
\end{equation}

Macro-F1 falls to 0.5526. In contrast, once the objective is active, performance remains comparatively stable across

\begin{equation}
\lambda_{NP}\in[0.10,1.00],
\end{equation}

with mean Macro-F1 ranging from 0.6489 to 0.6580.

The numerically highest value is obtained at \(\lambda_{NP}=0.50\), but the difference from the reference value 0.30 is approximately 0.0005, which is negligible relative to the observed seed-level variation. The result therefore does not support the claim that 0.30 is a uniquely optimal coefficient.

Instead, the more important finding is that the presence of targeted Normal--Phishing supervision matters substantially more than fine adjustment of its weight. This distinction is important for the interpretation of X-MoE: the most strongly supported mechanism is problem-specific supervision, whereas several architectural routing choices remain replaceable by simpler alternatives under the current dataset.


\subsection{Connectivity Determines the Value of Graph Evidence}

The connectivity-stratified analysis demonstrates that graph usefulness is conditional rather than uniform across addresses.

For connected nodes, removing isolate masking increases Macro-F1 from 0.6524 to 0.7273. The largest class-specific change occurs for Scamming, where F1 increases from 0.4363 to 0.6301. This result indicates that stronger access to structural evidence can improve fraud discrimination when usable neighbourhood information is present.

However, the same change introduces a small false-positive cost among isolated Normal nodes. The reference configuration A0 produces an isolated Normal false-positive rate of zero in the evaluated splits, whereas A4 increases this value to 0.000307, corresponding to approximately 3.1 additional false-positive predictions per 10,000 isolated Normal addresses.

The isolate-safe design should therefore be interpreted as a conservative modelling choice rather than as the configuration that maximises aggregate Macro-F1. It imposes the principle that graph-derived predictions should not influence an address for which no observed neighbourhood exists. This constraint protects isolated Normal addresses from unsupported structural contribution, but it also reduces the flexibility of the overall system.

The particularly strong Scamming improvement in A4 suggests that graph evidence is useful for at least part of the Scamming population represented in BLTE. In contrast, the corresponding Phishing improvement is comparatively small. This difference reinforces the importance of class-specific evaluation: Phishing and Scamming do not benefit from structural information to the same extent.

The present dataset does not allow this issue to be resolved completely. All labelled Phishing and Scamming test addresses in the evaluated splits are connected, whereas the isolated subset contains only Normal addresses. The analysis can therefore establish the false-positive protection associated with isolate masking for isolated Normal nodes, but it cannot determine whether isolated fraudulent addresses would benefit from or be harmed by the same mechanism.

The broader implication is that graph availability should be treated as part of the prediction condition rather than as a uniform property of the dataset. An address with informative neighbours and an isolated address do not provide equivalent evidence to a message-passing model, and reporting only aggregate graph-model performance can conceal this distinction.


\subsection{What STRICT-INDUCTIVE Reveals}

STRICT-INDUCTIVE provides a more demanding evaluation because validation and test nodes cannot exchange messages with other nodes in the same evaluation partition. Under this setting, X-MoE remains the strongest evaluated graph-dependent model, achieving a Macro-F1 of 0.5478 compared with 0.4401 for GraphSAGE.

The persistence of this advantage is important because it shows that the main comparison is not restricted to a topology-visible evaluation setting. X-MoE outperforms MLP and all three conventional GNN baselines across all ten matched strict-evaluation seeds, and all four comparisons remain statistically significant after Holm correction.

At the same time, X-MoE experiences the largest relative clean-to-strict decline, decreasing by 16.69\%, compared with 10.47\% for GraphSAGE. This reduction is informative rather than contradictory. X-MoE is explicitly designed to exploit relational evidence when it is available; restricting target-side graph visibility therefore removes part of the information the model is designed to use.

The decrease should not be interpreted as direct evidence that the topology-visible clean protocol contains label leakage. The two protocols differ in the amount of relational information available during message passing, and the observed performance change alone cannot establish the cause of the difference. The narrower and defensible conclusion is that part of the predictive advantage of X-MoE depends on graph evidence that becomes less accessible under stricter construction.

The class-specific strict results further support the earlier interpretation. The Phishing advantage remains pronounced, with X-MoE increasing F1 from 0.0825 for GraphSAGE to 0.3611. Scamming evidence is weaker when measured using threshold-dependent F1, although average precision remains improved.

STRICT-INDUCTIVE therefore strengthens the claim that X-MoE provides a comparative advantage within BLTE under restricted graph visibility, but it does not establish temporal, cross-chain, cross-dataset, or adversarial generalisation. Those settings involve different evidence conditions and remain outside the present evaluation.


\subsection{Implications for Graph Mixture-of-Experts Fraud Detection}

The combined results suggest several broader implications for the design and evaluation of graph mixture-of-experts models for blockchain fraud detection.

First, expert diversity should correspond to genuinely different evidence conditions. In X-MoE, the feature-oriented pathway can operate independently of graph neighbourhoods, while the GraphSAGE pathway can exploit relational evidence when connectivity exists. This distinction is more meaningful than introducing multiple experts that consume essentially identical information.

Second, specialised supervision may be more important than additional expert capacity under extreme class imbalance. The Normal--Phishing objective produces a much clearer empirical effect than the removal of several architectural routing components. For rare-class fraud detection, identifying and directly supervising the most consequential confusion boundary may therefore be more valuable than increasing model depth or routing complexity.

Third, routing complexity should be justified empirically. Shared routing and fixed outer fusion perform competitively, and in some cases exceed the full class-specific routing configuration. Future graph-MoE studies should therefore compare adaptive routing against simpler alternatives rather than assuming that a more expressive gating mechanism necessarily provides greater predictive value.

Fourth, connectivity-aware evaluation should accompany graph-based fraud detection whenever a substantial proportion of nodes lacks usable neighbourhood information. Aggregate performance alone cannot distinguish gains obtained from connected nodes from behaviour on isolated nodes. Connectivity-stratified reporting provides a clearer view of when relational modelling is genuinely contributing useful evidence.

Fifth, predictive quality and architectural complexity should be considered jointly. A2 retains almost the same Macro-F1 as the full model with substantially fewer parameters and lower observed training time. This result suggests that a reduced configuration may be preferable when computational efficiency is the dominant deployment requirement, whereas the full architecture may remain useful when explicit routing behaviour and evidence separation are important analytical objectives.

Overall, the empirical evidence supports a more nuanced interpretation of X-MoE than a simple claim of architectural superiority. The framework demonstrates that complementary feature and graph evidence, combined with targeted minority-class supervision, can improve balanced multi-class Ethereum fraud detection. At the same time, the experiments show that not every routing mechanism is necessary, graph contribution depends strongly on connectivity, and simpler configurations may provide favourable performance--complexity trade-offs. These findings define both the practical value and the empirical boundaries of the proposed framework.


\section{Limitations and Threats to Validity}

The empirical results support the effectiveness of X-MoE within the evaluated multi-class BLTE setting, but several limitations constrain the scope of the conclusions. These limitations concern dataset coverage, graph construction, generalisation, architectural dependence, statistical power, and computational interpretation.

\subsection{Dataset Scope and Label Completeness}

The evaluation is conducted using a single Ethereum dataset derived from the BLTE resource. This setting enables controlled comparison with existing multi-class Ethereum address-classification work, but the resulting conclusions remain specific to the characteristics of the evaluated data.

Only 69 addresses are labelled Phishing and 99 are labelled Scamming among 73,033 addresses. Although stratified splitting and ten matched random seeds reduce dependence on a single partition, the minority classes remain extremely small. Class-specific metrics are therefore sensitive to the small number of fraud-labelled examples assigned to each split.

The interpretation of the Normal class also requires caution. A Normal label denotes an address that is not labelled as Phishing or Scamming in the evaluated resource; it does not establish that the address has been independently verified as benign. Public fraud annotations may be incomplete, delayed, or biased toward detected incidents. The reported results should consequently be interpreted as supervised classification performance with respect to the available BLTE annotations rather than definitive adjudication of fraudulent or legitimate behaviour.

The study also considers a closed three-class setting. Addresses are assumed to belong to Normal, Phishing, or Scamming. Previously unseen fraud categories, mixed fraud behaviours, and open-set detection are outside the present formulation.


\subsection{Graph Construction and Connectivity}

The released graph contains 6,588 directed edges for 73,033 addresses, leaving approximately 92.32\% of nodes isolated. This sparsity is a property of the released representation used in the experiments and should not be interpreted as an intrinsic property of the complete Ethereum transaction network.

The experiments use the released processed graph rather than independently reconstructing transaction relationships from raw blockchain data. Missing edges may therefore reflect the upstream graph-construction process rather than the absence of interactions in Ethereum itself.

The topology-visible clean protocol retains the released graph across train, validation, and test partitions. Consequently, nodes assigned to different partitions may remain structurally connected even though feature preprocessing and labels are split correctly. This motivates the additional STRICT-INDUCTIVE evaluation, which removes target--target and cross-evaluation message passing.

STRICT-INDUCTIVE reduces this source of topology visibility but does not eliminate all limitations associated with evaluating nodes from the same underlying resource. The strict protocol should therefore be interpreted as a stronger node-inductive test rather than as an independent external validation dataset.

The connectivity-stratified analysis is also constrained by class composition. All labelled Phishing and Scamming test addresses in the evaluated splits are connected, whereas the isolated subset contains only Normal addresses. The experiment can therefore quantify false-positive behaviour among isolated Normal nodes but cannot determine the effectiveness of isolate-safe routing for isolated fraudulent addresses.


\subsection{Generalisation Beyond the Evaluated Setting}

The present experiments do not establish temporal generalisation. Training and test addresses are drawn from the same processed BLTE resource rather than from chronologically separated periods. Consequently, the evaluation does not measure robustness to future fraud campaigns, behavioural drift, newly created scam strategies, or changes in Ethereum transaction patterns.

The study also does not establish cross-dataset or cross-chain generalisation. Structural and behavioural signals observed in BLTE may differ from those available in independently collected Ethereum datasets or on other blockchain ecosystems such as Bitcoin, BNB Smart Chain, Solana, or Polygon.

STRICT-INDUCTIVE evaluates whether X-MoE retains its comparative advantage when target-side graph visibility is restricted, but this should not be interpreted as evidence of temporal or cross-domain transfer.

Adversarial robustness is likewise outside the current evaluation. The experiments do not test deliberate manipulation of transaction features, graph structure, connectivity, or routing behaviour by adaptive attackers.


\subsection{Architectural and Hyperparameter Dependence}

The reported results depend on a particular implementation of the Conflict-Aware Multitask Expert, independent GraphSAGE Expert, routing networks, auxiliary objectives, and optimisation procedure.

The ablation study reduces this concern by identifying components whose removal materially changes performance. However, it does not constitute an exhaustive architecture search. In particular, shared routing, fixed outer fusion, and removal of isolate masking can equal or exceed the reference X-MoE configuration in aggregate Macro-F1.

The complete architecture should therefore be interpreted as the reference framework evaluated in this study rather than as a globally optimal configuration.

The sensitivity analysis provides stronger evidence for the Normal--Phishing coefficient. Performance deteriorates substantially when the objective is removed but remains comparatively stable for

\begin{equation}
\lambda_{NP}\in[0.10,1.00].
\end{equation}

This result reduces concern that the principal finding depends on a narrowly selected value of \(\lambda_{NP}\), but other design parameters, including hidden dimensions, gate dimensions, dropout, and auxiliary-loss weights, were not subjected to exhaustive sensitivity analysis.

The study intentionally avoids selecting a revised architecture or coefficient retrospectively on the basis of the test-set sensitivity results.


\subsection{Statistical Validity}

Statistical comparisons are based on ten matched random seeds. The paired design controls for differences in data partitions and enables direct seed-level comparison, but the number of paired observations remains limited.

The clean-protocol comparison nevertheless shows consistent X-MoE improvement over MLP, GCN, GAT, and GraphSAGE across all ten seeds, with Holm-adjusted exact Wilcoxon \(p\)-values of 0.007812. These results provide evidence of consistency within the evaluated partitions but do not remove uncertainty associated with the small number of minority-class examples.

Holm correction is used to reduce the risk of false-positive findings across related comparison families. Non-significant results, however, should not be interpreted as proof of equivalence between configurations.

The 95\% confidence intervals and paired effect sizes reported in the study quantify uncertainty and practical effect magnitude under the evaluated seeds. They should not be interpreted as population-level guarantees for Ethereum fraud detection beyond the BLTE setting.


\subsection{Computational and Interpretability Limitations}

Observed training times depend on the specific CPU implementation and software environment used in this study. They are useful for comparing configurations executed under the same experimental conditions but should not be interpreted as universal runtime estimates.

The full X-MoE configuration contains more parameters and requires more training time than several simplified variants. In particular, the A2 configuration retains nearly the same Macro-F1 with substantially lower computational cost. The present experiments therefore do not establish that the complete model provides the optimal performance--efficiency trade-off.

The study also does not provide a comprehensive deployment evaluation. Inference latency, memory consumption under large-scale online processing, throughput, and performance under continuously arriving blockchain transactions are not evaluated.

Finally, X-MoE exposes routing weights, expert disagreement, connectivity indicators, and expert contributions that can support architectural inspection. These quantities provide diagnostic transparency but should not be interpreted as causal explanations of fraudulent behaviour. A large routing weight indicates greater reliance on a particular expert; it does not establish that a specific neighbour, transaction, or feature caused the predicted fraud label.

Taken together, these limitations bound the claims of the present study. The evidence supports improved multi-class BLTE classification through complementary expert modelling and targeted minority-class supervision, while broader conclusions concerning temporal robustness, external-dataset transfer, cross-chain generalisation, adversarial resistance, isolated-fraud detection, deployment scalability, and causal interpretability require further evaluation.



\section{Reproducibility and Open Implementation}

To support independent verification, the implementation of the evaluated models and experimental protocols is made publicly available at:

\begin{center}
\url{https://github.com/nurhalizaabdulwahab197/xmoe-ethereum-fraud-detection}
\end{center}

The repository contains the implementations of MLP, GCN, GAT, GraphSAGE, and X-MoE, together with deterministic experiment runners, reproduction instructions, dependency specifications, and reference outputs for the multi-class Ethereum fraud-detection experiments.

The released workflow reproduces the principal topology-visible clean protocol using ten random seeds,

\begin{equation}
\{100,101,\ldots,109\},
\end{equation}

300 training epochs, stratified 70/15/15 train/validation/test splitting, training-only feature standardisation, deterministic execution, and checkpoint selection based on the highest validation Macro-F1.

The same repository also documents the X-MoE ablation, connectivity-stratified evaluation, strict node-inductive protocol, and Normal--Phishing coefficient sensitivity analysis described in Sections V and VI.

The verified CPU software environment used for the reported experiments consists of Python 3.11.9, PyTorch 2.14.1+cpu, PyTorch Geometric 2.8.0.post1, NumPy 2.4.6, Pandas 3.0.6, scikit-learn 1.9.1, and SciPy 1.17.1.

Table~\ref{tab:reproducibility_environment} summarises the principal software dependencies.

\begin{table}[t]
\centering
\caption{Verified software environment used for the reported CPU experiments}
\label{tab:reproducibility_environment}

\renewcommand{\arraystretch}{1.12}
\setlength{\tabcolsep}{5pt}
\scriptsize

\begin{tabular}{ll}
\toprule
\textbf{Software} & \textbf{Version} \\
\midrule

Python & 3.11.9 \\
PyTorch & 2.14.1+cpu \\
PyTorch Geometric & 2.8.0.post1 \\
NumPy & 2.4.6 \\
Pandas & 3.0.6 \\
scikit-learn & 1.9.1 \\
SciPy & 1.17.1 \\

\bottomrule
\end{tabular}
\end{table}

To ensure that the intended processed BLTE artefacts are used, the repository records SHA-256 checksums for the two required files:

{\scriptsize
\begin{verbatim}
BLTE_multiclass_graph.pt
2859A8A3A815EDD02B5CD85D46C00D730B2B66C4E8D047145DB2E843C9E38456

BLTE_multiclass_v2.csv
02E376DDD98BA595372F19731FE5E51BBC90BF897B2A52B3C8BB98E786423B11
\end{verbatim}
}

These checksums are used to distinguish reproduction of the reported experiments from reconstruction of the original BLTE dataset. The present study operates on the processed artefacts released by the upstream project and does not independently regenerate the blockchain-data collection or labelling pipeline.

The strict node-inductive implementation further records graph-integrity diagnostics for each seed, including the numbers of training, validation, and test edges and explicit checks for forbidden direct validation--validation, test--test, and validation--test edges. Indirect paths through training nodes can remain in each admissible evaluation view. These diagnostics are intended to make the graph-visibility assumptions of STRICT-INDUCTIVE directly auditable.

The repository also includes the per-seed outputs used to compute aggregate means, standard deviations, class-specific metrics, statistical comparisons, and ablation results. The clean-protocol paired statistical analysis is based on the same ten matched partitions used for the reported MLP, GCN, GAT, GraphSAGE, and X-MoE comparisons.

In addition to the source implementation, an interactive demonstration of the X-MoE prediction workflow is available at:

\begin{center}
\url{https://xmoe-ethereum-fraud-detection.streamlit.app/}
\end{center}

The demonstration is intended to illustrate the prediction and model-interaction workflow. Quantitative claims in this paper are based on the controlled multi-seed experiments described in Sections V and VI rather than on individual demo predictions.


\section{Conclusion}

This study investigated whether task-aware expert fusion can improve multi-class Ethereum fraud detection under the combined conditions of extreme class imbalance and highly uneven graph connectivity.

To address this setting, we proposed X-MoE, a mixture-of-experts framework that combines a Conflict-Aware Multitask Expert with an independent GraphSAGE Expert. The framework separates feature-oriented and graph-structural evidence, applies class-specific routing conditioned on expert outputs, disagreement, and node connectivity, and incorporates targeted supervision for the difficult Normal--Phishing boundary.

Under the topology-visible clean protocol, X-MoE achieves a mean Macro-F1 of 0.6575, compared with 0.4916 for GraphSAGE, the strongest conventional graph baseline. X-MoE also outperforms the feature-only MLP, GCN, GAT, and GraphSAGE across all ten matched seeds, with all clean-protocol comparisons remaining significant after Holm correction.

The largest class-specific improvement occurs for Phishing. Relative to GraphSAGE, X-MoE increases Phishing F1 from 0.1023 to 0.5374 and average precision from 0.1871 to 0.5098. The ablation study provides the strongest mechanism-level evidence for the dedicated Normal--Phishing discrimination objective: removing this objective reduces overall Macro-F1 by 0.1049 and substantially degrades Phishing performance.

The experiments also show that architectural complexity is not uniformly beneficial. Shared routing, removal of isolate masking, and fixed outer mixing can match or exceed the reference model in aggregate Macro-F1, while a substantially smaller configuration retains nearly the same performance. Connectivity-stratified analysis further reveals a trade-off between stronger connected-node discrimination and conservative behaviour on isolated Normal addresses.

Under STRICT-INDUCTIVE, X-MoE retains the highest Macro-F1 among the evaluated graph-dependent models, achieving 0.5478 compared with 0.4401 for GraphSAGE and outperforming all three conventional GNN baselines across the ten matched seeds. The corresponding performance reduction relative to the clean protocol confirms that X-MoE remains sensitive to graph visibility and that part of its predictive advantage depends on relational evidence.

Sensitivity analysis further shows that the Normal--Phishing objective does not depend on a narrowly tuned coefficient. Performance deteriorates markedly when the objective is removed, but remains comparatively stable for positive values of \(\lambda_{NP}\) between 0.10 and 1.00.

Overall, the evidence supports X-MoE as a framework for improving balanced multi-class Ethereum fraud detection through complementary feature and graph evidence combined with targeted minority-class supervision. At the same time, the results demonstrate that graph utility depends on connectivity, routing complexity must be justified empirically, and simpler model variants can provide favourable performance--complexity trade-offs.

Future work should evaluate the framework on temporally separated data, independently collected Ethereum datasets, additional fraud categories, and other blockchain ecosystems. Further investigation is also required for isolated fraudulent addresses, open-set fraud recognition, adversarial robustness, online graph evolution, deployment efficiency, and external validation of the observed minority-class gains.



\balance
\bibliographystyle{IEEEtran}
\bibliography{references}

\end{document}
